\documentclass[11pt]{article}
\usepackage{amsmath,amssymb,amsthm,mathtools}
\usepackage{booktabs}
\usepackage[T1]{fontenc}
\usepackage{microtype}
\usepackage[hidelinks]{hyperref}
\usepackage[margin=1.15in]{geometry}

\newtheorem{theorem}{Theorem}[section]
\newtheorem{proposition}[theorem]{Proposition}
\newtheorem{lemma}[theorem]{Lemma}
\newtheorem{corollary}[theorem]{Corollary}
\newtheorem{remark}[theorem]{Remark}
\newtheorem{conjecture}[theorem]{Conjecture}

\newcommand{\F}{\mathbb F}
\newcommand{\C}{\mathbb C}
\newcommand{\Q}{\mathbb Q}
\newcommand{\Z}{\mathbb Z}
\newcommand{\supp}{\operatorname{supp}}
\newcommand{\rank}{\operatorname{rank}}

\title{Parity-Sensitive Fourier Uncertainty and Zero-Set Rigidity\\on the Finite Parabola}
\author{Dongwei Li\\
School of Mathematics, Hefei University of Technology\\
Hefei 230601, China\\
\texttt{dongweili@hfut.edu.cn}}
\date{}

\begin{document}
\maketitle

\begin{abstract}
Let $p\ge5$ be prime, let $R\subset\mathbb F_p$ satisfy
$1\le |R|\le p-1$, and set
\[
 F(a,b)=\sum_{x\in R}c_x\omega^{ax^2+bx},
 \qquad (a,b)\in\mathbb F_p^2,
 \qquad c_x\ne0,
\]
where $\omega=e^{2\pi i/p}$.  We prove a parity-sensitive uncertainty principle for complex Fourier spectra
supported on the finite parabola.
If $|R|=2r$, then
\[
 |Z(F)|\le p+2r-2,
\]
and the bound is sharp for every even support size $2\le |R|\le p-1$, including the endpoint $|R|=p-1$.  If $|R|=2r+1$, then
\[
 |Z(F)|\le \min\{p+2r-2,\,r(r+1)\}.
\]
Thus fixed odd spectral sparsity forces a number of zeros bounded
independently of $p$, whereas sufficiently large even zero sets are rigid:
for even support $|R|=2r$, if $|Z(F)|>r(r+1)$, then $F$ vanishes
identically on a nonvertical affine line and has at most $r(r-1)$ further
zeros.

The proof reduces arbitrary complex coefficients to cyclotomic data and
uses a truncated $(1-\omega)$-adic expansion.  The first nonzero finite-field
jets satisfy
\[
 \partial_b^2Q_j=\partial_aQ_{j-1},
\]
turning the two-dimensional zero problem into a multiplicity and
component-persistence problem for algebraic curves.

As an application, we solve the total-support uncertainty problem for the
standard complete set of $p+1$ mutually unbiased bases in $\mathbb C^p$:
\[
 \min_{0\ne\psi\in\mathbb C^p}
 \sum_j |\operatorname{supp}_{\mathcal B_j}(\psi)|
 =p^2-p+2.
\]
We also classify all extremizing projective states and give an exact
enumeration formula.
\end{abstract}

\medskip
\noindent\textbf{Keywords.}
uncertainty principles; finite Fourier analysis; finite parabola; zero-set rigidity;
mutually unbiased bases; extremizer classification.

\medskip
\noindent\textbf{2020 Mathematics Subject Classification.}
Primary 43A25; Secondary 11T23, 42C15, 81P55.

\section{Introduction}
\label{sec:introduction}

A basic theme in finite Fourier analysis is that a function and its Fourier
transform cannot both be sparse.  Classical group-theoretic and finite
versions include the results of Donoho--Stark, Smith and Meshulam; for a broad
modern survey of uncertainty principles and their variants, see
\cite{DonohoStark,SmithGroups,Meshulam,WigdersonSurvey}.  On the cyclic group
of prime order the phenomenon is exceptionally rigid: Chebotarev's
nonvanishing-minor theorem says that every square minor of the prime Fourier
matrix is nonzero; see, for example, the generalized Vandermonde treatment of
Evans--Isaacs and the recent extension of Loukaki
\cite{EvansIsaacs,LoukakiChebotarev}.  Tao used this rigidity to prove
\begin{equation}
 |\supp f|+|\supp\widehat f|\ge p+1
 \label{eq:intro-tao}
\end{equation}
for every nonzero $f:\F_p\to\C$ \cite{Tao}; see also the symmetry refinement
\cite{GarciaKaraaliKatz}.  The question considered here is what survives of
this rigidity in two dimensions when one Fourier support is constrained to a
curved one-dimensional set.  For the finite parabola we obtain a sharp
even-support bound, together with a striking parity dichotomy and a
classification of the large-zero-set mechanism.

For $R\subset\F_p$ and $c_x\ne0$, set
\begin{equation}
 F_{R,c}(a,b)
 =\sum_{x\in R}c_x\omega^{ax^2+bx},
 \qquad (a,b)\in\F_p^2,
 \label{eq:intro-parabolic-sum}
\end{equation}
where $\omega=e^{2\pi i/p}$.  Up to Fourier normalization and sign
conventions, this is the inverse Fourier transform of a function supported on
\[
 \Gamma=\{(x^2,x):x\in\F_p\}\subset\F_p^2.
\]
Write
\[
 Z(F_{R,c})=\{(a,b)\in\F_p^2:F_{R,c}(a,b)=0\}.
\]
Our first main theorem determines the zero-set scale and its parity-sensitive
structure.

\begin{theorem}[Parity-sensitive parabolic zero theorem]
\label{thm:intro-parabolic}
Let \(p\ge5\) be prime and let \(1\le |R|\le p-1\).

\begin{enumerate}
\item If \(|R|=1\), then \(F_{R,c}\) has no zeros.

\item If
\[
 |R|=2r,
\]
then
\[
 \boxed{|Z(F_{R,c})|\le p+2r-2.}
\]
The bound is sharp for every even support size \(2\le |R|\le p-1\).

\item If
\[
 |R|=2r+1,
\]
then
\[
 \boxed{
 |Z(F_{R,c})|
 \le
 \min\{p+2r-2,\ r(r+1)\}.
 }
\]

\item If \(|R|=2r\) and
\[
 |Z(F_{R,c})|>r(r+1),
\]
then there exists a nonvertical affine line \(\ell\subset\F_p^2\)
such that
\[
 F_{R,c}|_\ell\equiv0,
 \qquad
 |Z(F_{R,c})\setminus\ell|\le r(r-1).
\]
\end{enumerate}
\end{theorem}

The hypothesis in part~(4) is nonvacuous only when
\[
 r(r+1)<p+2r-2,
\]
since otherwise part~(2) already gives
$|Z(F_{R,c})|\le r(r+1)$.

The theorem shows that parity changes the order of magnitude of the extremal
zero set.  Even support permits line-driven configurations with $p+O(r)$
zeros, and the bound $p+2r-2$ is sharp throughout the stated range.  Odd
support admits no affine zero line and instead satisfies the uniform bound
$r(r+1)$, independent of the ambient prime.  Equivalently, if a nonzero
function on $\F_p^2$ has Fourier support of odd size $2r+1$ contained in the
parabola, then
\[
 |\supp f|\ge p^2-r(r+1).
\]
Thus fixed odd Fourier sparsity forces near-full physical support.  The large
even-support alternative is equally rigid: once the quadratic threshold
$r(r+1)$ is crossed, an entire affine line must appear and all remaining
zeros are quantitatively controlled.

A first consequence resolves the total-support problem for the standard
complete set of mutually unbiased bases in prime dimension.  Let
$\mathcal B_\infty$ be the computational basis and, for $a\in\F_p$, let
\[
 \phi_{a,b}(x)=p^{-1/2}\omega^{ax^2+bx},
 \qquad b,x\in\F_p,
\]
with $\mathcal B_a=\{\phi_{a,b}:b\in\F_p\}$.  These $p+1$ bases form the
standard complete MUB; see \cite{WF,BandyopadhyayMUB,KlappeneckerRottelerConstruction}
for foundational constructions and \cite{DurtMUB,McNultyWeigertReview} for
broader reviews.  Complete MUBs also carry a natural complex-projective
design structure \cite{KlappeneckerRottelerDesign,RoyScottDesign}.  For
$0\ne\psi\in\C^p$, write
\begin{equation}
 \mathcal S_p(\psi)
 :=\sum_{j\in\F_p\cup\{\infty\}}
 |\supp_{\mathcal B_j}(\psi)|,
 \qquad
 T_s(p):=\min_{\psi\ne0}\mathcal S_p(\psi).
 \label{eq:intro-Sp-Ts}
\end{equation}
Fiorentino and Weigert proved the general lower bound
\begin{equation}
 \mathcal S_p(\psi)\ge\frac{(p+1)^2}{2},
 \label{eq:intro-FW-bound}
\end{equation}
and obtained sharp information in the first small prime dimensions; their
work leaves the exact value of $T_s(p)$ open in general \cite{FW}.  Support
uncertainty and sparse representation for pairs of bases have a substantial
literature, including coherence-based formulations and the notions of
complete and higher-order incompatibility
\cite{EladBruckstein,DeBievre,DeBievre2023,XuIncompatibility}.
Those results concern two-basis transition geometry, whereas $T_s(p)$ is a
simultaneous extremal problem over the entire complete set of $p+1$ bases.
The parabolic theorem determines it for every prime $p\ge5$.

\begin{theorem}[Exact complete-MUB support uncertainty]
\label{thm:intro-exact}
For every prime $p\ge5$,
\[
 \boxed{T_s(p)=p^2-p+2.}
\]
\end{theorem}

Indeed, after moving a minimum-support MUB to the computational basis, the
coefficients in the remaining $p$ bases are encoded by a parabolic sum of the
form \eqref{eq:intro-parabolic-sum}.  If that minimum support has size $k$,
Theorem~\ref{thm:intro-parabolic} yields
\[
 \mathcal S_p(\psi)
 \ge
 p^2-p+k-2\Big\lfloor\frac{k}{2}\Big\rfloor+2,
\]
and hence the claimed lower bound; two-point states attain equality.

The equality theory is substantially more rigid than the minimum value
alone.  Every extremizer has even minimum support
\[
 2r\in\{2,4,\ldots,p-3\}.
\]
After normalizing a minimum-support MUB to the computational basis, its
support is a union of reflection pairs
\[
 \{c\pm\rho_1,\ldots,c\pm\rho_r\},
\]
and, after removal of a linear phase, the coefficients are antisymmetric
about $c$.  The remaining zeros are governed by a finite constraint matrix
whose maximal minors are all nonzero.  We prove a converse as well: every
admissible datum produces a unique extremizing projective ray.  This gives a
complete parametrization and a closed enumeration formula for all
extremizers.  The same theorem determines the minimum distance of the
associated complete-MUB analysis code and the cogirth of the represented
vector matroid.

Complete equality classifications have been obtained in other finite
uncertainty settings, for example for support inequalities on several finite
abelian groups \cite{BonamiGhobber} and for the short-time Fourier transform
on finite cyclic groups \cite{NicolaSTFT}.  Here the equality geometry is of
a different kind: it is generated by the curved spectral constraint and, in
the MUB formulation, must be compatible simultaneously across all $p+1$
bases.

We next describe the mechanism behind the parabolic theorem.  The argument
starts with arbitrary complex coefficients.  A minimal-support zero-row
relation has a one-dimensional kernel; maximal minors therefore place a
scaled coefficient vector in the cyclotomic field $\Q(\omega)$.  After
localization at $\pi=1-\omega$ and passage to a truncated jet algebra, one
obtains a formal function $\mathcal G$ satisfying
\begin{equation}
 \partial_b^2\mathcal G=T\,\partial_a\mathcal G.
 \label{eq:intro-schrodinger}
\end{equation}
If $Q_0,Q_1,\ldots$ are the jets beginning at the first nonzero order, then
\[
 \partial_b^2Q_0=0,
 \qquad
 \partial_b^2Q_j=\partial_aQ_{j-1},
\]
and therefore
\[
 Q_0(a,b)=A(a)+bB(a).
\]
This simple transverse geometry is the source of the sharp even bound.  A
nonexceptional vertical fiber contains at most one zero; if $A$ and $B$ have
a common root of multiplicity $m$, the recurrence gives at most $2m+1$ zeros
on the corresponding exceptional fiber.  Since the total common-root
multiplicity is bounded by $\deg B$, the two-dimensional zero count collapses
to a one-dimensional multiplicity budget.  The odd theorem requires a second
layer: common components of successive jets can persist only for controlled
orders, and a simultaneous odd/even induction converts that persistence
information into the uniform quadratic bound.

The result lies in a different regime from general finite-plane uncertainty.
Bir\'o and Lev obtain support inequalities for arbitrary functions on
$\F_p^2$, with bounds sensitive to the geometry of the two support sets
\cite{BiroLev}; Garcia, Karaali and Katz prove sharp refinements under group
symmetries \cite{GarciaKaraaliKatz}.  More recently, restriction estimates
have been used explicitly to strengthen uncertainty and annihilating-pair
inequalities on finite and locally compact abelian groups
\cite{IosevichMayeli,IosevichJamingMayeli}.  These results reinforce the role
of spectral geometry in uncertainty, but their conclusions are support- or
norm-based inequalities rather than exact cardinality bounds for the zero set.

Finite-field restriction theory was initiated by Mockenhaupt and Tao
\cite{MockenhauptTao} and has developed extensively for quadratic surfaces,
including paraboloids \cite{Lewko,LewkoKakeya,IKL}.  Recent work has also
begun a sharp theory of finite-field extension inequalities and their
extremizers, including the parabola and moment curves
\cite{GonzalezOliveira,BiswasEtAl}.  That theory optimizes extension norms;
our problem instead asks how many coefficients of the inverse transform can
vanish exactly when its spectrum is sparse and lies on the parabola.  The
answer exhibits a parity dichotomy not captured by the cited norm inequalities:
fixed odd sparsity gives a $p$-independent zero bound, while sufficiently large
even zero sets are forced by an affine zero line.

There is also a continuous analogue at the level of philosophy in the theory
of Heisenberg uniqueness pairs, where one studies Fourier transforms of
measures supported on curves that vanish on prescribed sets
\cite{HedenmalmMontes,GiriSrivastava}.  The present setting is different in
being finite and quantitative: the zero set itself is variable, and the goal
is its sharp maximal cardinality together with the structure of equality.
To our knowledge, the parity-sensitive parabolic zero theorem, its sharp
even-support range, and the accompanying large-zero-set rigidity do not
follow from the existing uncertainty, restriction, or uniqueness-pair
results.

\paragraph{Organization.}
Section~\ref{sec:setup} fixes notation.  Section~\ref{sec:parabolic-exact}
proves the parabolic zero theorem, and Section~\ref{sec:odd-refinements}
develops the finer odd-support results.  Section~\ref{sec:consequences}
derives the Fourier and MUB consequences.  Sections~\ref{sec:extremizer-structure}
and~\ref{sec:complete-extremizers} establish the equality classification and
enumeration.  Section~\ref{sec:code-matroid} records the coding and matroid
consequences, and Section~\ref{sec:discussion} discusses extensions and open
directions.

The logical dependence is simpler than the section order may suggest.  The
basic parabolic zero theorem of Section~\ref{sec:parabolic-exact} already
implies the exact complete-MUB value in Section~\ref{sec:consequences}; the
parity-sensitive component analysis of Section~\ref{sec:odd-refinements} is
not needed for that application.  It supplies instead the stronger standalone
odd-support bound and the large-even-zero-set rigidity, while
Sections~\ref{sec:extremizer-structure}--\ref{sec:complete-extremizers} analyze
equality in the basic bound and classify the MUB extremizers.

\section{Preliminaries}
\label{sec:setup}

Let $p\ge5$ be prime, let $\F_p$ be the field with $p$ elements, and put
\[
 \omega=e^{2\pi i/p}.
\]
For $a,b\in\F_p$, define
\begin{equation}
 \phi_{a,b}(x)
 =p^{-1/2}\omega^{ax^2+bx},
 \qquad x\in\F_p,
 \label{eq:standard-MUB}
\end{equation}
and let
\[
 \mathcal B_a=\{\phi_{a,b}:b\in\F_p\}.
\]
Together with the computational basis
$\mathcal B_\infty=\{e_x:x\in\F_p\}$, these form the standard complete set
of $p+1$ MUBs \cite{WF}.

For a vector $\psi\ne0$, write
\[
 \supp_{\mathcal B}(\psi)
 =\{v\in\mathcal B:\langle v,\psi\rangle\ne0\},
 \qquad
 s_{\mathcal B}(\psi):=|\supp_{\mathcal B}(\psi)|.
\]
Thus $s_{\mathcal B}(\psi)$ is a support cardinality, not a norm.
The quantities $\mathcal S_p(\psi)$ and $T_s(p)$ are defined in
\eqref{eq:intro-Sp-Ts}.  Since only vanishing matters, all statements are
projective in $\psi$.

For two distinct standard MUBs, the transition matrix is equivalent to the
prime Fourier matrix under multiplication of rows and columns by nonzero
phases and under row and column permutations.  To see this, use
Lemma~\ref{lem:standard-mub-transitivity} below to send one of the two bases
to \(\mathcal B_\infty\); the other is then some \(\mathcal B_a\).  The
transition matrix from \(\mathcal B_\infty\) to \(\mathcal B_a\) has entries
\[
 p^{-1/2}\omega^{ax^2+bx}
 =\omega^{ax^2}\bigl(p^{-1/2}\omega^{bx}\bigr),
\]
so it differs from the prime Fourier matrix only by a diagonal row phase.
Tao's theorem is invariant under these operations and therefore gives
\begin{equation}
 s_{\mathcal B_j}(\psi)
 +
 s_{\mathcal B_k}(\psi)
 \ge p+1
 \qquad(j\ne k).
 \label{eq:pairwise-MUB-support}
\end{equation}
This is the two-basis input behind the earlier complete-MUB bound
\eqref{eq:intro-FW-bound}; the stronger results below exploit the simultaneous
quadratic structure of all $p$ finite bases.

\begin{lemma}[Transitivity of the standard MUB]
\label{lem:standard-mub-transitivity}
For every basis \(\mathcal B_j\) in the standard complete set
\[
 \{\mathcal B_\infty\}\cup\{\mathcal B_a:a\in\F_p\},
\]
there exists a unitary \(U\) which permutes the whole standard complete
MUB, up to phases and permutations inside individual bases, and satisfies
\[
 U\mathcal B_j=\mathcal B_\infty.
\]
Consequently all support sizes are preserved, merely with the basis
labels permuted.
\end{lemma}

\begin{proof}
For \(c\in\F_p\), define the diagonal chirp
\[
 (C_cf)(x)=\omega^{cx^2}f(x).
\]
From the definition
\[
 \phi_{a,b}(x)=p^{-1/2}\omega^{ax^2+bx}
\]
one has
\[
 C_c\mathcal B_\infty=\mathcal B_\infty,
 \qquad
 C_c\mathcal B_a=\mathcal B_{a+c}.
\]
Let \(\mathsf F\) denote the unitary Fourier matrix
\[
 \mathsf F_{x,y}=p^{-1/2}\omega^{xy}.
\]
Then \(\mathsf F^*\mathcal B_0=\mathcal B_\infty\) and
\(\mathsf F^*\mathcal B_\infty=\mathcal B_0\), up to permutations of the basis
vectors.  If \(a\ne0\), then
\[
 (\mathsf F^*\phi_{a,b})(y)
 =
 \frac1p\sum_{x\in\F_p}\omega^{ax^2+(b-y)x}.
\]
Completing the square gives
\[
 ax^2+(b-y)x
 =
 a\left(x+\frac{b-y}{2a}\right)^2
 -
 \frac{(b-y)^2}{4a}.
\]
The quadratic Gauss sum
\[
 G(a):=\sum_{x\in\F_p}\omega^{ax^2}
\]
is nonzero.  Indeed, since \(p\) is odd, the change of variables
\[
 u=x-y,\qquad v=x+y
\]
is a bijection of \(\F_p^2\), and therefore
\[
 \begin{aligned}
 |G(a)|^2
 &=\sum_{x,y\in\F_p}\omega^{a(x^2-y^2)}
 =\sum_{u,v\in\F_p}\omega^{auv}\\
 &=\sum_{u\in\F_p}
   \left(\sum_{v\in\F_p}\omega^{auv}\right)
 =p.
 \end{aligned}
\]
Thus \(|G(a)|=\sqrt p\).  Hence
\[
 (\mathsf F^*\phi_{a,b})(y)
 =
 \frac{G(a)}p
 \omega^{-(b-y)^2/(4a)}
 =
 \zeta_{a,b}\,
 p^{-1/2}\omega^{-y^2/(4a)+(b/(2a))y},
\]
where \(\zeta_{a,b}\) is a nonzero scalar of modulus one.  Therefore
\[
 \mathsf F^*\mathcal B_a=\mathcal B_{-1/(4a)}
 \qquad(a\ne0),
\]
up to phases and a permutation of the second label.  Thus \(\mathsf F^*\) also
permutes the standard complete MUB.

If \(j=\infty\), take \(U=I\).  If \(j=a\in\F_p\), first apply
\(C_{-a}\), which sends \(\mathcal B_a\) to \(\mathcal B_0\), and then
apply \(\mathsf F^*\), which sends \(\mathcal B_0\) to
\(\mathcal B_\infty\).  The composition still permutes the complete
standard MUB.
\end{proof}


\section{Sharp zero counts for parabolic Fourier sums}
\label{sec:parabolic-exact}

Throughout this section \(p\ge5\) is prime and
\(\omega=e^{2\pi i/p}\).  If
\[
 \psi=\sum_{x\in R}c_xe_x
\]
has computational support \(R\subset\F_p\), then under the standard Hermitian
inner-product convention, for physical finite-MUB labels \(A,\beta\in\F_p\),
\begin{equation}
 \langle \phi_{A,\beta},\psi\rangle
 =p^{-1/2}\sum_{x\in R}c_x\omega^{-A x^2-\beta x}
 =p^{-1/2}F_{R,c}(-A,-\beta).
 \label{eq:mub-parabola-label-map}
\end{equation}
where
\begin{equation}
 F_{R,c}(a,b)
 :=
 \sum_{x\in R}c_x\omega^{a x^2+b x},
 \qquad (a,b)\in\F_p^2.
 \label{eq:parabolic-sum-rigorous}
\end{equation}
Since \((A,\beta)\mapsto(-A,-\beta)\) is a bijection of \(\F_p^2\), the
number of zero coefficients among the \(p\) finite quadratic bases is exactly
\(\#Z(F_{R,c})\).
For such a sum write
\[
 Z(F_{R,c})
 :=\{(a,b)\in\F_p^2:F_{R,c}(a,b)=0\}.
\]
Whenever an auxiliary coefficient vector is allowed to have zero coordinates,
its \emph{actual support} means the set of indices on which its coefficients
are nonzero.  Thus every parabolic sum is always interpreted on its actual
support when a support-size theorem is applied.
The proof below establishes a quantitative upper bound for
$|Z(F_{R,c})|$.


\paragraph{Proof architecture.}
There are four steps.  First, a rank induction reduces arbitrary complex
coefficients either to smaller support or to a one-dimensional kernel defined
over \(K=\Q(\omega)\).  Second, cyclotomic localization converts exact zeros
into vanishing conditions for truncated jets over \(\F_p\).  Third, the first
nonzero jet is shown by a Vandermonde moment argument to occur by order
\(\lfloor k/2\rfloor\).  Finally, the parabolic jet recurrence reduces the
two-dimensional zero count to a multiplicity budget for the pencil
\(A(a)+bB(a)\).

\subsection{Algebraic preliminaries}

\begin{lemma}[Truncated cyclotomic model]
\label{lem:truncated-cyclotomic-model}
Let
\[
 K=\mathbb Q(\omega),\qquad
 \pi=1-\omega,
\]
and let
\[
 \mathcal O:=\mathbb Z[\omega]_{(\pi)}
\]
be the localization at the prime ideal generated by \(\pi\).  Then
\begin{equation}
 \mathcal O/(\pi^{p-1})
 \cong
 \F_p[X]/(X^{p-1}),
 \qquad
 \pi\longmapsto X.
 \label{eq:cyclotomic-truncation}
\end{equation}
Moreover the substitution
\begin{equation}
 X=1-e^T,
 \qquad
 e^T:=\sum_{j=0}^{p-2}\frac{T^j}{j!},
 \label{eq:formal-change-rigorous}
\end{equation}
defines an isomorphism
\[
 \F_p[X]/(X^{p-1})
 \cong
 \F_p[T]/(T^{p-1}).
\]
Under the composite isomorphism,
\begin{equation}
 \omega^h=(1-\pi)^h
 \longmapsto
 e^{hT}
 \pmod{T^{p-1}}
 \label{eq:omega-to-exp}
\end{equation}
for every \(h\in\F_p\).
\end{lemma}

\begin{proof}
The identity
\[
 \Phi_p(1)=\prod_{j=1}^{p-1}(1-\omega^j)=p
\]
shows that \(p\) is associated in \(\mathcal O\) to
\(\pi^{p-1}\): indeed
\[
 1-\omega^j
 =
 (1-\omega)(1+\omega+\cdots+\omega^{j-1}),
\]
and the second factor reduces to \(j\ne0\) modulo \(\pi\), hence is a
unit in \(\mathcal O\).  Thus
\[
 p=u\pi^{p-1}
\]
for some \(u\in\mathcal O^\times\).  In particular
\(\mathcal O/(\pi^{p-1})\) has characteristic \(p\).

Since \(\omega\equiv1\pmod\pi\), reduction modulo \(\pi\) sends
\(\mathbb Z[\omega]\) onto \(\mathbb Z/(p)\); hence the residue field is
\[
 \mathcal O/(\pi)\cong\F_p.
\]
The ring \(\mathbb Z[\omega]\) is the ring of integers of \(K\)
\cite[Ch.~2]{Washington}, hence is Dedekind, so its localization
\(\mathcal O\) at the nonzero prime \((\pi)\) is a discrete valuation ring
with uniformizer \(\pi\).
Successive reduction modulo powers of \(\pi\) shows that every class
modulo \(\pi^{p-1}\) can be written
\[
 a_0+a_1\pi+\cdots+a_{p-2}\pi^{p-2},
 \qquad a_j\in\F_p.
\]
Hence the homomorphism
\[
 \F_p[X]/(X^{p-1})
 \longrightarrow
 \mathcal O/(\pi^{p-1}),
 \qquad X\longmapsto\pi,
\]
is surjective.  Both rings have \(p^{p-1}\) elements: the left side has
\(\F_p\)-dimension \(p-1\), while the right side has a filtration
\[
 \mathcal O/(\pi^{p-1})
 \supset
 (\pi)/(\pi^{p-1})
 \supset\cdots\supset
 (\pi^{p-2})/(\pi^{p-1})
 \supset0
\]
whose \(p-1\) successive quotients are all isomorphic to the residue
field \(\F_p\).  The surjection is therefore an isomorphism, proving
\eqref{eq:cyclotomic-truncation}.

Because \(1,\ldots,p-2\) are invertible in \(\F_p\), the truncated
exponential in \eqref{eq:formal-change-rigorous} is well defined.
Furthermore
\[
 1-e^T=-T+O(T^2)
\]
has zero constant term and nonzero linear coefficient.  A substitution
of this form is an automorphism of the local Artinian algebra
\(\F_p[T]/(T^{p-1})\).  Explicitly, its inverse is given by the
truncated logarithm
\[
 T=\log(1-X)
 =
 -\sum_{j=1}^{p-2}\frac{X^j}{j}
 \pmod{X^{p-1}}.
\]
Finally, choose the standard integer representative
\(h\in\{0,1,\ldots,p-1\}\).  The formal identity
\[
 (e^T)^h=e^{hT}
\]
is first understood in \(\mathbb Q[[T]]\).  Reduce this identity modulo
\(T^{p-1}\).  Only the coefficients of degrees \(0,\ldots,p-2\) remain,
and every denominator occurring there divides a product of integers strictly
smaller than \(p\).  Those denominators are therefore units in
\(\mathbb Z_{(p)}\), so the truncated identity lies in
\(\mathbb Z_{(p)}[T]/(T^{p-1})\) and may be reduced coefficientwise modulo
\(p\).  Since \(1-X=e^T\) in the truncated algebra, this gives
\[
 (1-X)^h=e^{hT}\pmod{T^{p-1}},
\]
which proves \eqref{eq:omega-to-exp}.
\end{proof}

\begin{lemma}[Reduction from complex to cyclotomic coefficients]
\label{lem:complex-to-cyclotomic}
Fix \(R\subset\F_p\), \(|R|=k\ge2\), and a set
\(Z\subset\F_p^2\).  For \(y=(a,b)\in Z\), put
\[
 v_y:=
 \bigl(\omega^{a x^2+b x}\bigr)_{x\in R}
 \in K^R.
\]
Suppose there exists \(0\ne c\in\C^R\) such that
\[
 \sum_{x\in R}c_x(v_y)_x=0
 \qquad(y\in Z).
 \label{eq:zero-linear-system}
\]
Let \(r_Z\) be the rank over \(K\) of the matrix whose rows are the
\(v_y\).

Then exactly one of the following alternatives is available:

\begin{enumerate}
\item \(r_Z\le k-2\), in which case there exists
\(0\ne c'\in\C^R\) satisfying all equations
\eqref{eq:zero-linear-system} and having at most \(r_Z+1\) nonzero
coordinates (and hence at most \(k-1\));
\item \(r_Z=k-1\), in which case the common kernel is one-dimensional
and is generated by a vector in \(K^R\).  Hence \(c\) is a nonzero
complex scalar multiple of a vector in \(K^R\).
\end{enumerate}
\end{lemma}

\begin{proof}
Because the coefficient matrix has entries in \(K\), its rank is
unchanged after extending scalars from \(K\) to \(\C\).  Since the
nonzero vector \(c\) belongs to its kernel,
\[
 r_Z\le k-1.
\]

Assume first \(r_Z\le k-2\).  Let
\[
 W:=
 \left\{
 u\in\C^R:
 \sum_{x\in R}u_x(v_y)_x=0
 \text{ for every }y\in Z
 \right\}.
\]
Then
\[
 \dim_\C W=k-r_Z.
\]
Choose any set \(S\subset R\) of exactly \(k-r_Z-1\) coordinates and let
\[
 H_S:=\{u\in\C^R:u_x=0\text{ for every }x\in S\}.
\]
The subspace \(H_S\) has codimension \(k-r_Z-1\).  Therefore
\[
 \dim(W\cap H_S)
 \ge (k-r_Z)-(k-r_Z-1)=1.
\]
Choose \(0\ne c'\in W\cap H_S\).  It satisfies every equation
\eqref{eq:zero-linear-system} and can be nonzero on at most
\[
 k-|S|=r_Z+1
\]
coordinates.  This proves the first alternative.

Now assume \(r_Z=k-1\).  Choose \(k-1\) linearly independent rows.  The
resulting \((k-1)\times k\) matrix \(M\) has a one-dimensional kernel.
The vector of signed maximal minors
\[
 w_j=(-1)^{j-1}\det M^{(j)},
\]
where \(M^{(j)}\) is obtained by deleting the \(j\)-th column, is a
nonzero vector in \(K^R\) and satisfies \(Mw=0\).  Therefore the kernel
over \(K\) is \(Kw\), while after extending scalars the kernel over
\(\C\) is \(\C w\).  Hence \(c=\lambda w\) for some
\(\lambda\in\C^\times\).
\end{proof}

\subsection{Cyclotomic jets}

\begin{lemma}[First nonzero jet]
\label{lem:first-nonzero-jet}
Let
\[
 R\subset\F_p,\qquad 2\le k:=|R|\le p-1,
\]
and let \(c_x\in K^\times\) for every \(x\in R\).  After multiplication by a common
nonzero scalar, assume
\begin{equation}
 c_x\in\mathcal O
 \quad(x\in R),
 \qquad
 \min_{x\in R}v_\pi(c_x)=0.
 \label{eq:primitive-c-rigorous}
\end{equation}
Map each \(c_x\) through Lemma~\ref{lem:truncated-cyclotomic-model} and
write its image as
\[
 C_x(T)\in\F_p[T]/(T^{p-1}).
\]
Define
\begin{equation}
 \mathcal G(T;a,b)
 :=
 \sum_{x\in R}
 C_x(T)e^{T(a x^2+b x)}
 =
 \sum_{j=0}^{p-2}H_j(a,b)T^j
 \pmod{T^{p-1}}.
 \label{eq:G-rigorous}
\end{equation}
Then:

\begin{enumerate}
\item \(H_j\in\F_p[a,b]\) has total degree at most \(j\);
\item if \(n\) is the least index for which \(H_n\) is not the zero
polynomial, then
\begin{equation}
 n\le\left\lfloor\frac{k}{2}\right\rfloor.
 \label{eq:n-rigorous}
\end{equation}
\end{enumerate}
\end{lemma}

\begin{proof}
Write
\[
 C_x(T)=\sum_{\ell=0}^{p-2}c_{x,\ell}T^\ell.
\]
The coefficient of \(T^j\) in
\[
 C_x(T)e^{T(a x^2+b x)}
\]
is
\[
 \sum_{\ell=0}^{j}
 c_{x,\ell}
 \frac{(a x^2+b x)^{j-\ell}}{(j-\ell)!}.
\]
Hence \(H_j\) has total degree at most \(j\), proving the first assertion.

Put
\[
 r:=\left\lfloor\frac{k}{2}\right\rfloor.
\]
Suppose, for contradiction, that
\[
 H_0=H_1=\cdots=H_r=0
\]
as polynomials.  For each \(0\le j\le r\), the homogeneous component of
total degree \(j\) of \(H_j\) comes only from the terms with
\(\ell=0\), and therefore equals
\begin{equation}
 \frac1{j!}
 \sum_{x\in R}
 c_{x,0}(a x^2+b x)^j.
 \label{eq:top-homogeneous-rigorous}
\end{equation}
Expanding,
\[
 (a x^2+b x)^j
 =
 \sum_{t=0}^{j}
 \binom jt
 a^t b^{j-t}x^{j+t}.
\]
Since \(j<p\), all coefficients \(j!\) and \(\binom jt\) occurring here
are nonzero in \(\F_p\).  Vanishing of
\eqref{eq:top-homogeneous-rigorous} therefore implies
\begin{equation}
 \sum_{x\in R}c_{x,0}x^d=0
 \qquad
 (j\le d\le2j).
 \label{eq:moment-window}
\end{equation}
As \(j\) runs from \(0\) to \(r\), the intervals
\([j,2j]\) cover every integer
\[
 0,1,\ldots,k-1.
\]
Thus
\[
 \sum_{x\in R}c_{x,0}x^d=0
 \qquad(0\le d\le k-1).
\]
Fix an ordering
\[
 R=\{x_1,\ldots,x_k\}.
\]
The Vandermonde matrix
\[
 V=(x_j^d)_{\substack{0\le d\le k-1\\1\le j\le k}}
\]
has determinant
\[
 \det V=\prod_{1\le i<j\le k}(x_j-x_i)\ne0
\]
in \(\F_p\), because the elements of \(R\) are distinct.  Hence
\(c_{x,0}=0\) for every \(x\in R\).

On the other hand, the primitivity condition
\eqref{eq:primitive-c-rigorous} means that at least one \(c_x\) is a unit
of \(\mathcal O\).  Its image modulo \(\pi\), which is precisely
\(c_{x,0}\), is nonzero.  This contradiction proves
\eqref{eq:n-rigorous}.
\end{proof}

\begin{lemma}[Parabolic jet recurrence]
\label{lem:schrodinger-jet}
Retain the notation of Lemma~\ref{lem:first-nonzero-jet}, let \(n\) be
the first nonzero jet order, and put
\[
 Q_j(a,b):=H_{n+j}(a,b)
\]
whenever \(n+j\le p-2\).  Then
\begin{equation}
 \partial_b^2Q_0=0,
 \qquad
 \partial_b^2Q_j=\partial_aQ_{j-1}
 \quad(j\ge1).
 \label{eq:schrodinger-recurrence-rigorous}
\end{equation}
Consequently
\begin{equation}
 Q_0(a,b)=A(a)+bB(a)
 \label{eq:leading-pencil-rigorous}
\end{equation}
for polynomials \(A,B\in\F_p[a]\) satisfying
\begin{equation}
 \deg A\le n,
 \qquad
 \deg B\le n-1.
 \label{eq:AB-degree-rigorous}
\end{equation}
Moreover, whenever \(n+j\le p-2\),
\begin{equation}
 \deg_bQ_j\le2j+1.
 \label{eq:b-degree-rigorous}
\end{equation}
\end{lemma}

\begin{proof}
In the truncated polynomial ring,
\[
 \partial_b e^{T(a x^2+b x)}
 =
 Tx\,e^{T(a x^2+b x)}
\]
and
\[
 \partial_a e^{T(a x^2+b x)}
 =
 Tx^2\,e^{T(a x^2+b x)}.
\]
Therefore
\begin{equation}
 \partial_b^2\mathcal G
 =
 T\,\partial_a\mathcal G
 \pmod{T^{p-1}}.
 \label{eq:formal-schrodinger-rigorous}
\end{equation}
Comparing coefficients of \(T^n\) gives
\(\partial_b^2H_n=0\).  Comparing coefficients of
\(T^{n+j}\), \(j\ge1\), gives
\[
 \partial_b^2H_{n+j}
 =
 \partial_aH_{n+j-1}.
\]
This is \eqref{eq:schrodinger-recurrence-rigorous}.

Since \(\deg Q_0\le n<p\), write
\[
 Q_0(a,b)=\sum_{d=0}^{n}q_d(a)b^d.
\]
The equation \(\partial_b^2Q_0=0\) gives
\[
 d(d-1)q_d(a)=0
 \qquad(d\ge2).
\]
Because \(2\le d<p\), the scalar \(d(d-1)\) is nonzero in \(\F_p\).
Thus \(q_d=0\) for every \(d\ge2\), proving
\eqref{eq:leading-pencil-rigorous}.  Since the total degree of \(Q_0\)
is at most \(n\), the degree bounds
\eqref{eq:AB-degree-rigorous} follow.

It remains to prove \eqref{eq:b-degree-rigorous}.  We use induction on
\(j\).  The case \(j=0\) has just been proved.  Suppose
\(\deg_bQ_{j-1}\le2j-1\).  Then
\[
 \deg_b(\partial_aQ_{j-1})\le2j-1.
\]
Also \(n+j\le p-2\), hence every power of \(b\) occurring in \(Q_j\)
has exponent \(<p\).  If \(Q_j\) contained a nonzero term
\(q_d(a)b^d\) with \(d\ge2j+2\), then the term
\[
 d(d-1)q_d(a)b^{d-2}
\]
would survive in \(\partial_b^2Q_j\), because \(d<p\), and would have
\(b\)-degree at least \(2j\).  This contradicts
\[
 \partial_b^2Q_j=\partial_aQ_{j-1}.
\]
Therefore \(\deg_bQ_j\le2j+1\).
\end{proof}

\begin{remark}[Algebraic support and jet annihilators]
\label{rem:general-jet-annihilator}
Identity~\eqref{eq:formal-schrodinger-rigorous} is the quadratic
instance of a more general mechanism.  If a frequency set
\(\Lambda\subset\F_p^m\) lies in \(P(\xi)=0\), where
\[
 P(z)=\sum_\alpha p_\alpha z^\alpha,\qquad D=\deg P,
\]
then the cyclotomic generating function is annihilated by
\[
 \mathcal D_{P,T}
 =
 \sum_\alpha
 p_\alpha T^{D-|\alpha|}\partial_y^\alpha.
\]
At the first nonzero jet only the highest homogeneous part survives:
\[
 P_D(\partial_y)Q_0=0.
\]
For the parabola \(u-v^2=0\), this becomes
\(\partial_b^2Q_0=0\), which is exactly the transverse linearity
\eqref{eq:leading-pencil-rigorous}.  The broad
Fourier--differential correspondence is classical; the role here is
its finite cyclotomic first-jet form.
\end{remark}

\begin{lemma}[Exact zeros vanish in every available jet]
\label{lem:exact-zero-all-jets}
Under the hypotheses of Lemma~\ref{lem:first-nonzero-jet}, define
\[
 F(a,b)=\sum_{x\in R}c_x\omega^{a x^2+b x}.
\]
If
\[
 F(\alpha,\beta)=0
\]
in \(K\), then
\begin{equation}
 H_j(\alpha,\beta)=0
 \qquad(0\le j\le p-2).
 \label{eq:exact-zero-all-jets}
\end{equation}
Equivalently, if \(Q_j\) is defined, then
\[
 Q_j(\alpha,\beta)=0.
\]
\end{lemma}

\begin{proof}
By \eqref{eq:primitive-c-rigorous}, every summand of \(F(\alpha,\beta)\)
lies in \(\mathcal O\).  The exact equality
\[
 F(\alpha,\beta)=0
\]
therefore remains zero after reduction modulo \(\pi^{p-1}\).
Under the isomorphism of Lemma~\ref{lem:truncated-cyclotomic-model}, this
reduced element is exactly
\[
 \mathcal G(T;\alpha,\beta)
 =
 \sum_{j=0}^{p-2}H_j(\alpha,\beta)T^j
 \in\F_p[T]/(T^{p-1}).
\]
The residue classes \(1,T,\ldots,T^{p-2}\) are linearly independent
over \(\F_p\).  Hence every coefficient is zero, which is precisely
\eqref{eq:exact-zero-all-jets}.
\end{proof}

\subsection{The zero bound}

\begin{lemma}[Cyclotomic parabolic zero bound]
\label{lem:cyclotomic-zero-rigorous}
Let
\[
 R\subset\F_p,\qquad 2\le k:=|R|\le p-1,
\]
and let \(c_x\in K^\times\) for every \(x\in R\).  Define
\[
 F(a,b)=\sum_{x\in R}c_x\omega^{a x^2+b x}.
\]
Then
\begin{equation}
 \#Z(F)
 :=
 \#\{(a,b)\in\F_p^2:F(a,b)=0\}
 \le
 p+2\left\lfloor\frac{k}{2}\right\rfloor-2.
 \label{eq:cyclotomic-zero-rigorous}
\end{equation}
\end{lemma}

\begin{proof}
Scale \(c\) so that \eqref{eq:primitive-c-rigorous} holds, and let
\(n\) be the first nonzero jet order.  By
Lemma~\ref{lem:first-nonzero-jet},
\begin{equation}
 n\le r:=\left\lfloor\frac{k}{2}\right\rfloor.
 \label{eq:n-vs-r}
\end{equation}
Write
\[
 Q_0(a,b)=A(a)+bB(a)
\]
as in Lemma~\ref{lem:schrodinger-jet}.

\medskip
\noindent
\textbf{Case 1: \(B\not\equiv0\).}
Let
\[
 D=\gcd(A,B),
\]
chosen monic.  We use the convention
\[
 \operatorname{ord}_\alpha 0=+\infty.
\]
Fix \(\alpha\in\F_p\).

If \(B(\alpha)\ne0\), then
\[
 A(\alpha)+bB(\alpha)=0
\]
has exactly one solution \(b\in\F_p\).  Since every exact zero must
annihilate \(Q_0\) by Lemma~\ref{lem:exact-zero-all-jets}, the vertical
fiber \(a=\alpha\) contains at most one exact zero.

If \(B(\alpha)=0\) but \(A(\alpha)\ne0\), then the fiber contains no
exact zeros.

It remains to consider a common root
\(\alpha\) of \(A\) and \(B\).  Put
\[
 m_\alpha
 :=
 \min\{\operatorname{ord}_\alpha A,
        \operatorname{ord}_\alpha B\}.
\]
Because \(B\not\equiv0\),
\[
 1\le m_\alpha\le\deg B\le n-1.
\]
Therefore
\[
 n+m_\alpha
 \le2n-1
 \le2r-1
 \le p-2,
\]
so \(Q_{m_\alpha}\) is defined.

Iterating the recurrence
\eqref{eq:schrodinger-recurrence-rigorous} \(m_\alpha\) times gives
\begin{equation}
 \partial_b^{2m_\alpha}Q_{m_\alpha}
 =
 \partial_a^{m_\alpha}Q_0.
 \label{eq:iterated-recurrence-rigorous}
\end{equation}
Indeed, the statement follows by induction:
\[
 \partial_b^{2s}Q_{m_\alpha}
 =
 \partial_a^sQ_{m_\alpha-s}
 \qquad(0\le s\le m_\alpha).
\]

At \(a=\alpha\),
\[
 \partial_a^{m_\alpha}Q_0(\alpha,b)
 =
 A^{(m_\alpha)}(\alpha)
 +
 bB^{(m_\alpha)}(\alpha).
 \label{eq:nonzero-linear-derivative}
\]
By definition of \(m_\alpha\), both \(A\) and \(B\) vanish to order at
least \(m_\alpha\), and at least one vanishes to exactly that order.
Since \(m_\alpha<p\), the factor \(m_\alpha!\) is nonzero in \(\F_p\).
Hence the polynomial in
\eqref{eq:nonzero-linear-derivative} is not identically zero.

It follows from \eqref{eq:iterated-recurrence-rigorous} that
\(Q_{m_\alpha}(\alpha,b)\) is not the zero polynomial.  By
\eqref{eq:b-degree-rigorous},
\[
 \deg_bQ_{m_\alpha}(\alpha,b)\le2m_\alpha+1<p.
\]
Every exact zero on this fiber is a root of this polynomial, by
Lemma~\ref{lem:exact-zero-all-jets}.  Hence the fiber contains at most
\begin{equation}
 2m_\alpha+1
 \label{eq:fiber-bound-common-root}
\end{equation}
exact zeros.

Let \(Z_{\F_p}(B)\) be the set of roots of \(B\) in \(\F_p\), and let
\(Z_{\F_p}(D)\) be the common roots of \(A\) and \(B\) in \(\F_p\).
Adding the fiber bounds gives
\[
 \begin{aligned}
 \#Z(F)
 &\le
 p-|Z_{\F_p}(B)|
 +\sum_{\alpha\in Z_{\F_p}(D)}(2m_\alpha+1)\\
 &=
 p
 +2\sum_{\alpha\in Z_{\F_p}(D)}m_\alpha
 +|Z_{\F_p}(D)|-|Z_{\F_p}(B)|\\
 &\le
 p+2\sum_{\alpha\in Z_{\F_p}(D)}m_\alpha\\
 &\le
 p+2\deg D\\
 &\le
 p+2\deg B\\
 &\le
 p+2n-2.
 \end{aligned}
 \label{eq:count-case-one-rigorous}
\]

\medskip
\noindent
\textbf{Case 2: \(B\equiv0\).}
Then \(Q_0=A(a)\) with \(A\ne0\).  By
Lemma~\ref{lem:exact-zero-all-jets}, an exact zero can occur only on a
fiber \(a=\alpha\) for which \(A(\alpha)=0\).

Let \(\alpha\in\F_p\) be a root of multiplicity
\[
 m_\alpha:=\operatorname{ord}_\alpha A.
\]
If
\[
 n+m_\alpha\le p-2,
 \label{eq:safe-root-order}
\]
then the same recurrence argument as in Case 1 gives
\[
 \partial_b^{2m_\alpha}Q_{m_\alpha}(\alpha,b)
 =
 A^{(m_\alpha)}(\alpha)\ne0.
\]
Thus \(Q_{m_\alpha}(\alpha,b)\) is a nonzero polynomial of degree at
most \(2m_\alpha+1\), and the fiber contains at most
\(2m_\alpha+1\) exact zeros.

Suppose first that \eqref{eq:safe-root-order} holds for every
\(\F_p\)-root of \(A\).  Then
\[
 \begin{aligned}
 \#Z(F)
 &\le
 \sum_{\alpha\in Z_{\F_p}(A)}(2m_\alpha+1)\\
 &=
 2\sum_{\alpha\in Z_{\F_p}(A)}m_\alpha
 +|Z_{\F_p}(A)|\\
 &\le
 2\deg A+\deg A\\
 &\le3n.
 \end{aligned}
\]
For \(n\ge1\),
\[
 3n\le p+2n-2
\]
because \(n\le r\le(p-1)/2\le p-2\).  If \(n=0\), then \(A\) is a
nonzero constant and \(\#Z(F)=0\).

It remains to treat a root for which
\[
 n+m_\alpha\ge p-1.
\]
Since
\[
 m_\alpha\le\deg A\le n
 \quad\text{and}\quad
 n\le r\le\frac{p-1}{2},
\]
all inequalities must be equalities:
\[
 k=p-1,
 \qquad
 n=m_\alpha=\frac{p-1}{2}.
 \label{eq:boundary-equalities}
\]
In particular \(\deg A=n\), and \(A\) has the single root
\(\alpha\) with multiplicity \(n\).  Hence every exact zero of \(F\)
lies on the single fiber \(a=\alpha\).

On that fiber,
\[
 b\longmapsto F(\alpha,b)
 =
 \sum_{x\in R}
 \bigl(c_x\omega^{\alpha x^2}\bigr)\omega^{bx}
\]
is the ordinary Fourier transform on \(\F_p\) of the nonzero function
\[
 x\longmapsto c_x\omega^{\alpha x^2}\mathbf 1_R(x),
\]
whose support is \(R\), of size \(p-1\).  Tao's
prime-dimensional uncertainty principle \cite{Tao} gives
\[
 |\supp F(\alpha,\cdot)|
 \ge
 p+1-(p-1)=2.
\]
Thus this fiber contains at most \(p-2\) zeros.  Since all other fibers
contain none,
\[
 \#Z(F)\le p-2.
\]
This is stronger than the desired bound.

Combining both cases,
\[
 \#Z(F)\le p+2n-2
 \le
 p+2\left\lfloor\frac{k}{2}\right\rfloor-2,
\]
which proves \eqref{eq:cyclotomic-zero-rigorous}.
\end{proof}

\begin{theorem}[Parabolic zero theorem]
\label{thm:parabolic-zero-rigorous}
Let
\[
 R\subset\F_p,\qquad 1\le k:=|R|\le p-1,
\]
and let \(c_x\in\C^\times\) for every \(x\in R\).  Define
\[
 F(a,b)=\sum_{x\in R}c_x\omega^{a x^2+b x}.
\]
If \(k=1\), then \(F\) has no zeros.  If \(k\ge2\), then
\begin{equation}
 \boxed{
 \#Z(F)
 \le
 p+2\left\lfloor\frac{k}{2}\right\rfloor-2.
 }
 \label{eq:parabolic-zero-rigorous}
\end{equation}
Equivalently,
\[
 \#Z(F)\le
 \begin{cases}
 p+k-2,&k\text{ even},\\
 p+k-3,&k\ge3\text{ odd}.
 \end{cases}
\]
\end{theorem}

\begin{proof}
We use induction on \(k\).

For \(k=1\),
\[
 F(a,b)=c_x\omega^{a x^2+b x}
\]
for the unique \(x\in R\), so \(F(a,b)\ne0\) everywhere.

Assume \(k\ge2\), and assume the theorem has been proved for every
strictly smaller support size.  Let
\[
 Z:=Z(F).
\]
For \(y=(a,b)\in Z\), define
\[
 v_y=
 \bigl(\omega^{a x^2+b x}\bigr)_{x\in R}.
\]
The coefficient vector \(c\) satisfies the linear system
\[
 \sum_{x\in R}c_x(v_y)_x=0
 \qquad(y\in Z).
\]
Apply Lemma~\ref{lem:complex-to-cyclotomic}.

If the rank \(r_Z\le k-2\), that lemma provides a nonzero vector \(c'\)
satisfying all the same zero equations and having at most \(k-1\) nonzero
coordinates.  Let
\[
 R':=\{x\in R:c'_x\ne0\},
\]
so \(1\le |R'|<k\), and let \(F'\) be the corresponding parabolic sum
written on its actual support \(R'\).
Then
\[
 Z(F)\subseteq Z(F').
\]
If \(|R'|=1\), then \(F'\) has no zeros, so the required estimate is
immediate.  If \(|R'|\ge2\), the induction hypothesis gives
\[
 \begin{aligned}
 |Z(F)|
 &\le |Z(F')|\\
 &\le
 p+2\left\lfloor\frac{|R'|}{2}\right\rfloor-2\\
 &\le
 p+2\left\lfloor\frac{k}{2}\right\rfloor-2.
 \end{aligned}
\]

If \(r_Z=k-1\), Lemma~\ref{lem:complex-to-cyclotomic} shows that
\(c\) is a nonzero complex scalar multiple of a vector in \(K^R\).
Multiplication by a nonzero scalar does not change the zero set.
Therefore Lemma~\ref{lem:cyclotomic-zero-rigorous} applies and gives
the required estimate.

These two cases exhaust all possibilities because the existence of the
nonzero kernel vector \(c\) implies \(r_Z\le k-1\).
\end{proof}

\section{Parity-sensitive refinements for odd support}
\label{sec:odd-refinements}

Throughout this section, unless otherwise stated, $p\ge5$ is prime and the actual support satisfies $1\le |R|\le p-1$.

The general parabolic zero theorem is optimal throughout the even-support range needed for the complete-MUB problem, but odd support is substantially more rigid.  The goal of this section is to prove the full parity-sensitive theorem announced in the introduction.  We isolate here the additional mechanisms: affine-line rigidity, low-support B\'ezout arguments, persistence and breaking of jet components, and restriction to rational components.  The component analysis below closes all full-rank branches.  Coupled with an even-support large-zero-line theorem, it yields a general parity-sensitive bound for every odd support size.  None of this section is needed for Theorem~\ref{thm:exact-mub-rigorous}.

\subsection{Affine lines and the first odd supports}

\begin{lemma}[Restriction to an affine line]
\label{lem:affine-line-zero-bound}
Let
\[
 F(a,b)=\sum_{x\in R}c_x\omega^{a x^2+b x},\qquad |R|=k,
\]
with all $c_x\ne0$, and let $\ell\subset\F_p^2$ be an affine line.
If $\ell$ is vertical, then
\[
 |Z(F)\cap\ell|\le k-1.
\]
If $\ell$ is nonvertical, then either $F|_\ell\equiv0$ or
\begin{equation}
 |Z(F)\cap\ell|
 \le
 \min\left\{k-1,\frac{p-1}{2}\right\}.
 \label{eq:nonvertical-line-sharp}
\end{equation}
The identically-zero alternative is impossible when $k$ is odd.  Hence
\begin{equation}
 |Z(F)\cap\ell|\le k-1
 \qquad(k\text{ odd}),
 \label{eq:odd-line-bound}
\end{equation}
and for odd support on a nonvertical line one may use the sharper
bound \eqref{eq:nonvertical-line-sharp}.
\end{lemma}

\begin{proof}
First consider the vertical line \(a=\alpha\).  Define
\[
 f_\alpha(x):=c_x\omega^{\alpha x^2}\mathbf 1_R(x).
\]
This is a nonzero function on \(\F_p\) with support exactly \(R\), hence
\(|\supp f_\alpha|=k\).  Up to the harmless sign convention in the Fourier
exponent, its Fourier transform is
\[
 b\longmapsto F(\alpha,b).
\]
Tao's prime-order uncertainty principle gives
\[
 k+|\supp F(\alpha,\cdot)|\ge p+1.
\]
Therefore \(|\supp F(\alpha,\cdot)|\ge p+1-k\), so the vertical fiber
contains at most \(p-(p+1-k)=k-1\) zeros.

Now let the line be nonvertical and write it as \(b=\mu a+\nu\).  Put
\[
 d_x:=c_x\omega^{\nu x},
 \qquad
 h(x):=x^2+\mu x,
\]
and define a function \(g:\F_p\to\C\) by
\[
 g(t):=\sum_{\substack{x\in R\\ h(x)=t}}d_x.
\]
Grouping the terms of the parabolic sum by the value of \(h(x)\) gives
\[
 F(a,\mu a+\nu)
 =\sum_{x\in R}d_x\omega^{a h(x)}
 =\sum_{t\in\F_p}g(t)\omega^{at}.
\]
If \(g\ne0\), Tao's theorem applied to \(g\) gives
\[
 |\supp g|+|\supp \widehat g|\ge p+1.
\]
The number of zeros of \(a\mapsto F(a,\mu a+\nu)\) is therefore at most
\(|\supp g|-1\).  Since
\[
 h(x)=\left(x+\frac\mu2\right)^2-\frac{\mu^2}{4},
\]
its image has cardinality \((p+1)/2\), and of course
\(|\supp g|\le |R|=k\).  Thus
\[
 |\supp g|\le\min\left\{k,\frac{p+1}{2}\right\},
\]
which proves \eqref{eq:nonvertical-line-sharp}.

It remains to analyze the alternative \(g=0\), which is exactly the condition
that \(F\) vanish identically on the line.  The fibers of \(h\) are the
orbits of the involution
\[
 x\longmapsto-\mu-x.
\]
There is one fixed point \(-\mu/2\); all other fibers have two points.  Since
every coefficient \(d_x\) with \(x\in R\) is nonzero, an occupied fiber
cannot meet \(R\) in exactly one point, for then the corresponding value of
\(g\) would be nonzero.  The fixed point therefore does not belong to \(R\),
and every occupied fiber contributes exactly two support points.  Hence
\(|R|\) is even.  In particular, exact vanishing on a nonvertical line is
impossible for odd support.
\end{proof}

\begin{lemma}[Affine symmetries of the parabolic zero problem]
\label{lem:parabolic-affine-symmetry}
Let \(t,\nu\in\F_p\).  If
\[
 R_t:=R-t,
 \qquad
 c^{(t)}_y:=c_{y+t}\quad(y\in R_t),
\]
then
\begin{equation}
 F_{R,c}(a,b)
 =
 \omega^{at^2+bt}
 F_{R_t,c^{(t)}}(a,b+2at).
 \label{eq:support-translation-symmetry}
\end{equation}
Moreover, replacing \(c_x\) by \(c_x\omega^{\nu x}\) replaces
\(F(a,b)\) by \(F(a,b+\nu)\), while replacing \(c_x\) by
\(c_x\omega^{\sigma x^2}\) replaces it by \(F(a+\sigma,b)\).  Hence support
translation, linear rephasing, and quadratic rephasing preserve the number
of zeros and preserve vertical/nonvertical affine-line structure.  The first
two operations act transitively on nonvertical affine lines, so any such line
may be normalized to \(b=0\); quadratic rephasing translates the vertical
coordinate \(a\).

If the coefficients have been cyclotomically normalized, these operations
preserve the first nonzero jet order and the total degree of its leading jet.
They also preserve divisibility of the first two jets by the transformed
nonvertical factor: if \(C\mid Q_0,Q_1\), then the corresponding transformed
factor divides the transformed \(Q_0,Q_1\).
\end{lemma}

\begin{proof}
Writing \(x=y+t\) gives \eqref{eq:support-translation-symmetry} by
expanding
\[
 a(y+t)^2+b(y+t)
 =ay^2+(b+2at)y+(at^2+bt).
\]
Linear and quadratic rephasing give the identities
\[
 \sum_x c_x\omega^{\nu x}\omega^{ax^2+bx}=F(a,b+\nu),
 \qquad
 \sum_x c_x\omega^{\sigma x^2}\omega^{ax^2+bx}=F(a+\sigma,b).
\]
In the truncated cyclotomic model, support translation gives
\[
 \mathcal G_{R,c}(T;a,b)
 =e^{T(at^2+bt)}
  \mathcal G_{R_t,c^{(t)}}(T;a,b+2at).
\]
Write the unit factor as \(1+T\Lambda+O(T^2)\).  Since
\(H_0=\cdots=H_{n-1}=0\), the transformed first two jets have the form
\[
 \widetilde Q_0=Q_0\circ\Theta,
 \qquad
 \widetilde Q_1=Q_1\circ\Theta+\Lambda\,(Q_0\circ\Theta),
\]
where \(\Theta\) is the corresponding affine change of \((a,b)\).  Thus a
common factor of \(Q_0\) and \(Q_1\) is carried to a common factor of the
transformed first two jets.  The same formulas show that the first nonzero
\(T\)-order is unchanged.  The affine substitutions preserve total degree.
Linear and quadratic rephasing are simply the substitutions
\(b\mapsto b+\nu\) and \(a\mapsto a+\sigma\), respectively.
\end{proof}

\begin{remark}[Why the parity mechanism characterizes quadratics]
\label{rem:segre-quadratic-characterization}
The preceding pairing argument has a finite-geometric converse.
For a general graph spectrum
\[
 \Gamma_\phi=\{(\phi(x),x):x\in\F_p\},
\]
exact vanishing on a nonvertical line is equivalent to fiberwise
cancellation for \(x\mapsto\phi(x)+\mu x\).  Thus odd support is
universally incapable of producing an exact zero line if every such
fiber has size at most two.  Conversely, if one fiber contains three
distinct points \(x_1,x_2,x_3\), taking coefficients \(1,1,-2\) on
those points gives an odd-support sum that vanishes identically on the
corresponding nonvertical line.  Hence universal odd-line rigidity is
equivalent to every fiber having size at most two, or equivalently to
\(\Gamma_\phi\) having no three collinear points.

Add the common point \(P_\infty=[1:0:0]\) at infinity of the vertical
lines.  The resulting \(p+1\) points form an oval in \(PG(2,p)\), so
Segre's theorem~\cite{SegreOval} implies, for odd \(p\), that they lie
on a nondegenerate conic.  Write its homogeneous equation as
\[
 AU^2+BUV+CV^2+DUW+EVW+FW^2=0.
\]
Since \(P_\infty\) lies on the conic, \(A=0\).  On the vertical affine
line \(V=xW\), the equation becomes
\[
 U(Bx+D)+Cx^2+Ex+F=0.
\]
For every \(x\in\F_p\), the affine graph point on this vertical
line lies on the conic.  If \(Bx+D=0\), the remaining constant term must
also vanish, so the whole vertical line would be a component of the conic,
contradicting nondegeneracy.  Hence \(Bx+D\ne0\) for every
\(x\in\F_p\).  Therefore \(B=0\) and \(D\ne0\), and nondegeneracy
forces \(C\ne0\).  Consequently
\[
 \phi(x)=-\frac{C}{D}x^2-\frac{E}{D}x-\frac{F}{D},
\]
a nondegenerate quadratic.  The converse is immediate because a
nonzero quadratic polynomial has fibers of size at most two.  Thus the
parity phenomenon used below is genuinely quadratic, rather than a
generic feature of polynomial graph spectra.
\end{remark}

Put
\[
 \mu_p:=\{\omega^j:j\in\F_p\},
\]
the group of \(p\)-th roots of unity.

\begin{theorem}[Exact three-point bound]
\label{thm:three-point-exact}
If $|R|=3$, then
\begin{equation}
 \boxed{|Z(F)|\le2,}
 \label{eq:three-point-bound}
\end{equation}
and the constant $2$ is attained.
\end{theorem}

\begin{proof}
Write $R=\{x_0,x_1,x_2\}$ and factor the $x_0$ phase.  The linear map
\[
 (a,b)\longmapsto
 \bigl(a(x_1^2-x_0^2)+b(x_1-x_0),\,
       a(x_2^2-x_0^2)+b(x_2-x_0)\bigr)
\]
has determinant $(x_1-x_0)(x_2-x_0)(x_1-x_2)\ne0$.  Thus zeros are in bijection with pairs $(z,w)\in\mu_p^2$ satisfying
\[
 c_0+c_1z+c_2w=0.
\]
For a solution, $|c_0+c_1z|=|c_2|$.  Hence $c_1z$ lies on two nonconcentric Euclidean circles, so there are at most two possibilities for $z$.  Equality occurs, for example, at $(z,w)=(1,1)$ and $(\omega,\omega^{-1})$ with $c_2=1$, $c_1=\omega^{-1}$, $c_0=-(1+\omega^{-1})$.
\end{proof}

We shall repeatedly use the following standard interpolation identity.

\begin{lemma}[Barycentric moments]
\label{lem:barycentric-moments}
For distinct $R=\{x_1,\ldots,x_k\}\subset\F_p$, $2\le k<p$, put
\[
 P_R(X)=\prod_{i=1}^k(X-x_i),\qquad w_i=P_R'(x_i)^{-1}.
\]
Then
\[
 \sum_iw_ix_i^m=0\quad(0\le m\le k-2),\qquad
 \sum_iw_ix_i^{k-1}=1,
\]
and
\[
 \sum_iw_ix_i^k=\sum_ix_i.
\]
Moreover the solutions of $\sum_iu_ix_i^m=0$ for $0\le m\le k-3$ are exactly
\[
 u_i=w_i(\lambda x_i+\mu).
\]
\end{lemma}

\begin{proof}
For \(1\le i\le k\), let
\[
 L_i(X):=\frac{P_R(X)}{(X-x_i)P_R'(x_i)}.
\]
Then \(L_i(x_j)=\delta_{ij}\), and the coefficient of \(X^{k-1}\) in
\(L_i\) is \(w_i=P_R'(x_i)^{-1}\).  If \(0\le m\le k-1\), Lagrange
interpolation gives
\[
 X^m=\sum_{i=1}^k x_i^mL_i(X).
\]
Comparing coefficients of \(X^{k-1}\) yields
\[
 \sum_iw_ix_i^m=0\quad(0\le m\le k-2),
 \qquad
 \sum_iw_ix_i^{k-1}=1.
\]

For \(m=k\), write
\[
 P_R(X)=X^k-e_1X^{k-1}+\text{terms of degree at most }k-2,
 \qquad
 e_1=\sum_i x_i.
\]
Modulo \(P_R\), the polynomial \(X^k\) is therefore congruent to a
polynomial of degree at most \(k-1\) whose \(X^{k-1}\)-coefficient is
\(e_1\).  This remainder has the same values \(x_i^k\) at the interpolation
nodes.  Applying the preceding coefficient comparison to the remainder gives
\[
 \sum_iw_ix_i^k=e_1=\sum_i x_i.
\]

Finally consider the \((k-2)\times k\) moment matrix
\[
 M=(x_i^m)_{0\le m\le k-3,\ 1\le i\le k}.
\]
It has rank \(k-2\), because any \((k-2)\times(k-2)\) minor obtained from
distinct columns is a nonzero Vandermonde determinant.  Hence \(\ker M\) has
dimension two.  For any \(\lambda,\mu\in\F_p\), the vector
\[
 u_i=w_i(\lambda x_i+\mu)
\]
lies in this kernel, since for \(0\le m\le k-3\)
\[
 \sum_i u_ix_i^m
 =\lambda\sum_iw_ix_i^{m+1}+\mu\sum_iw_ix_i^m=0.
\]
The two parameters give a two-dimensional family: if
\(\lambda x_i+\mu=0\) for every \(i\), two distinct nodes force
\(\lambda=\mu=0\).  Therefore this family is the whole kernel.
\end{proof}


\begin{lemma}[Maximal first-jet order forces full degree]
\label{lem:maximal-order-full-degree}
Let the actual support have size $2\le k\le p-1$, and let
\[
 n=\left\lfloor\frac{k}{2}\right\rfloor
\]
be the first nonzero cyclotomic jet order.  Then
\begin{equation}
 \boxed{\deg Q_0=n.}
 \label{eq:maximal-order-full-degree}
\end{equation}
\end{lemma}

\begin{proof}
The identities
\[
 H_0=\cdots=H_{n-1}=0
\]
give, from their highest homogeneous parts,
\[
 S_m:=\sum_{x\in R}c_{x,0}x^m=0
 \qquad(0\le m\le2n-2).
\]
If \(\deg Q_0<n\), then the highest homogeneous part of
\(H_n=Q_0\) also vanishes, giving
\[
 S_m=0
 \qquad(n\le m\le2n).
\]
Thus
\[
 S_0=\cdots=S_{2n}=0.
\]
If \(k=2n+1\), these are exactly \(k\) consecutive moments; if
\(k=2n\), the first \(k\) of them already suffice.  In either case
the Vandermonde matrix on the distinct support points forces every
\(c_{x,0}=0\), contradicting primitivity.
\end{proof}

\begin{lemma}[A full-degree leading jet cannot persist]
\label{lem:Q0-not-divide-Q1-universal}
Let \(n\ge2\) be the first nonzero cyclotomic jet, write
\[
 Q_0=A(a)+bB(a),
\]
and assume
\[
 \deg Q_0=n.
\]
Then
\begin{equation}
 \boxed{Q_0\nmid Q_1.}
 \label{eq:Q0-no-wholesale-persistence}
\end{equation}
\end{lemma}

\begin{proof}
If \(Q_0\mid Q_1\), then
\[
 \deg(Q_1/Q_0)\le(n+1)-n=1.
\]
Thus
\[
 Q_1=Q_0(c_0+c_1a+db).
\]
The recurrence \(\partial_b^2Q_1=\partial_aQ_0\) gives
\[
 2dB=A'+bB'.
\]
Thus $B'=0$ and $A'$ is constant.  Since all degrees are below $p$, $B$ is constant and $\deg A\le1$, contradicting $\deg Q_0\ge2$.
\end{proof}


\paragraph{Canonical nonvertical factor.}
Whenever a leading pencil is written as \(Q_0=A(a)+bB(a)\) with
\(B\not\equiv0\), we use the canonical factorization
\[
 D=\gcd(A,B),
 \qquad
 Q_0=D(a)C(a,b).
\]
Then \(C=A_0(a)+bB_0(a)\) with \(\gcd(A_0,B_0)=1\) and
\(B_0\ne0\).  The coprimality remains true over
\(\overline{\F}_p[a]\): a common root over \(\overline{\F}_p\) would
produce a nonconstant common factor over \(\F_p\) by taking the product of
its Frobenius conjugates.  Hence \(C\), viewed as a polynomial of degree one
in \(b\) over \(\overline{\F}_p[a]\), is primitive and cannot have a
nontrivial factorization.  Thus \(C\) is absolutely irreducible.  We use
this fact when applying B\'ezout over \(\overline{\F}_p\).

\paragraph{A B\'ezout refinement used below.}
The following elementary projective correction will be used repeatedly in the low-support and component arguments.

\begin{lemma}[B\'ezout with forced intersection at infinity]
\label{lem:bezout-infinity}
Let coprime $P,Q\in\overline{\F}_p[a,b]$ have degrees $D,E$ and $b$-degrees at most $s<D$ and $t<E$.  Then the number of common affine zeros satisfies
\begin{equation}
 \#Z_{\mathrm{aff}}(P,Q)
 \le DE-(D-s)(E-t).
 \label{eq:bezout-infinity}
\end{equation}
\end{lemma}

\begin{proof}
Let \(P^h(A,B,Z)\) and \(Q^h(A,B,Z)\) be the homogenizations of \(P\) and
\(Q\).  Because the \(b\)-degrees satisfy \(s<D\) and \(t<E\), both
projective curves pass through
\[
 P_\infty=[0:1:0].
\]
Work in the affine chart \(B=1\) around this point, with local coordinates
\[
 u=A/B,
 \qquad
 v=Z/B.
\]
A monomial \(a^ib^j\) of \(P\), after homogenization to total degree \(D\),
becomes \(A^iB^jZ^{D-i-j}\), hence locally becomes
\[
 u^iv^{D-i-j}.
\]
Its local total degree is \(D-j\ge D-s\).  Therefore the multiplicity of
\(P^h=0\) at \(P_\infty\) is at least \(D-s\).  The same argument gives
multiplicity at least \(E-t\) for \(Q^h=0\).

Since \(P\) and \(Q\) are coprime over \(\overline{\F}_p\), their
homogenizations have no common projective component.  The local intersection
inequality gives
\[
 I_{P_\infty}(P^h,Q^h)\ge(D-s)(E-t)
\]
\cite[Sec.~3.3]{FultonCurves}.  Projective B\'ezout gives total intersection
multiplicity \(DE\).  Every distinct affine common zero contributes local
intersection multiplicity at least one, so the number of affine common zeros
is at most
\[
 DE-I_{P_\infty}(P^h,Q^h)
 \le DE-(D-s)(E-t).
\]
\end{proof}


\begin{proposition}[Four-point threshold]
\label{prop:four-point-large-zero-structure}
If $|R|=4$ and $|Z(F)|>6$, then there is a nonvertical affine line $\ell$ such that
\[
 F|_\ell\equiv0,
 \qquad
 |Z(F)\setminus\ell|\le2.
\]
\end{proposition}

\begin{proof}
Let $r_Z$ be the zero-row rank.  If $r_Z\le2$, a coordinate-hyperplane cut gives a nonzero vector on at most three support points preserving all zeros.  Theorem~\ref{thm:three-point-exact} forces actual support two, whose zero set is one affine line; more than three zeros exclude a vertical line by Lemma~\ref{lem:affine-line-zero-bound}, and hence the original four-point restriction vanishes identically there.

Assume $r_Z=3$.  After cyclotomic reduction the first jet order satisfies $n\le2$.  The cases $n=0,1$ are immediate from $Q_0$.  For $n=2$, use Lemma~\ref{lem:parabolic-affine-symmetry} to translate $R$ so that $\sum_{x\in R}x=0$.  From $H_0=H_1=0$ and Lemma~\ref{lem:barycentric-moments}, the primitive residue vector is $\bar c_{x_i}=\lambda w_i$, whence
\[
 Q_0^{[2]}=\lambda ab.
\]
Thus $\deg Q_0=2$, and Lemma~\ref{lem:Q0-not-divide-Q1-universal} gives $Q_0\nmid Q_1$.  If $Q_0,Q_1$ are coprime, B\'ezout gives at most six zeros.  Otherwise they share exactly one linear factor.  If it were not Frobenius-stable, its distinct conjugate would also divide both polynomials, forcing the quadratic $Q_0$ to divide $Q_1$; hence the common line is defined over $\F_p$.  After factoring it out, the residual degrees are $1$ and at most $2$, so there are at most two zeros off the line.  More than six total zeros then give at least five on the line, which by Lemma~\ref{lem:affine-line-zero-bound} forces exact vanishing there.
\end{proof}

\begin{theorem}[Uniform five-point bound]
\label{thm:five-point-six-zero}
For every prime $p\ge7$,
\begin{equation}
 |R|=5\qquad\Longrightarrow\qquad\boxed{|Z(F)|\le6.}
 \label{eq:five-point-six}
\end{equation}
\end{theorem}

\begin{proof}
Assume $|Z(F)|\ge7$.  If $r_Z\le2$, two coordinate cuts reduce to support at most three, contradicting Theorem~\ref{thm:three-point-exact} and the odd line bound.  If $r_Z=3$, one cut gives support at most four; Proposition~\ref{prop:four-point-large-zero-structure} then places at least five original zeros on one affine line, again contradicting \eqref{eq:odd-line-bound}.

Let $r_Z=4$.  After cyclotomic reduction, $n\le2$.  The cases $n=0,1$ are immediate.  For $n=2$, use Lemma~\ref{lem:parabolic-affine-symmetry} to translate so that $\sum R=0$.  The residue moment equations through degree $2$ and Lemma~\ref{lem:barycentric-moments} give
\[
 \bar c_{x_i}=w_i(ux_i+v),\qquad (u,v)\ne(0,0),
\]
and hence
\[
 Q_0^{[2]}=uab+\frac v2a^2\ne0.
\]
Therefore $\deg Q_0=2$, and Lemma~\ref{lem:Q0-not-divide-Q1-universal} applies.  Coprime $Q_0,Q_1$ have at most six common zeros by B\'ezout.  Otherwise they share a linear factor \(L\) over
\(\overline{\F}_p\).  Because both \(Q_0\) and \(Q_1\) have coefficients
in \(\F_p\), every Frobenius conjugate of \(L\) is also a common factor.
If \(L\) were not defined over \(\F_p\), its distinct conjugate would give
two different linear factors of the quadratic \(Q_0\); their product would
therefore be a nonzero scalar multiple of \(Q_0\), forcing \(Q_0\mid Q_1\),
contrary to Lemma~\ref{lem:Q0-not-divide-Q1-universal}.  Hence the common
line is defined over \(\F_p\).  After removing it, the remaining factors
have degrees one and at most two and are coprime, so B\'ezout gives at most
two common points off the line.  The original five-point sum cannot vanish
identically on a nonvertical line because its support size is odd; by
\eqref{eq:odd-line-bound} it has at most four exact zeros on the common
line.  Thus \(|Z(F)|\le4+2=6\).
\end{proof}


\begin{proposition}[Five-point sharpness]
\label{prop:five-point-sharpness}
The constant in Theorem~\ref{thm:five-point-six-zero} is sharp.  In fact,
for $p=11$ there is a five-point parabolic sum with exactly six zeros.
\end{proposition}

\begin{proof}
Let $\zeta=e^{2\pi i/11}$, let
\[
 R=\{0,1,2,3,4\},
\]
and consider the four points
\[
 y_1=(1,4),\qquad y_2=(3,0),\qquad
 y_3=(6,0),\qquad y_4=(7,1).
\]
Form the $4\times5$ cyclotomic evaluation matrix
\[
 A=\bigl(\zeta^{a_i x^2+b_i x}\bigr)_
 {\substack{1\le i\le4\\0\le x\le4}},
 \qquad y_i=(a_i,b_i),
\]
and let $A^{(x)}$ denote the matrix obtained by deleting the column indexed
by $x$.  The signed maximal-minor vector
\[
 c_x=\zeta^{10}(-1)^x\det A^{(x)},\qquad 0\le x\le4,
\]
lies in $\ker A$.  Reducing modulo
\[
 \Phi_{11}(X)=1+X+\cdots+X^{10}
\]
gives the following representatives:
\begin{align*}
 c_0&=\zeta^9-3\zeta^6+2\zeta^5-\zeta+1,\\
 c_1&=\zeta^9+\zeta^8-\zeta^5+3\zeta^4+\zeta^3
       +2\zeta^2+2\zeta+2,\\
 c_2&=3\zeta^9-2\zeta^8+\zeta^7-\zeta^6-\zeta^3+2\zeta-2,\\
 c_3&=3\zeta^7+\zeta^5+\zeta^4+2\zeta^3+3\zeta^2+2\zeta-1,\\
 c_4&=\zeta^8-2\zeta^6+2\zeta^5-3\zeta^3+2\zeta^2-\zeta+1.
\end{align*}
Each is nonzero, since it is represented by a nonzero polynomial of degree
at most nine whereas $\Phi_{11}$ has degree ten.  Hence the actual support
is all five points of $R$.  The cofactor identity gives
\[
 F(y_i)=0\qquad(1\le i\le4).
\]
For the two further points one obtains, by direct polynomial reduction,
\begin{align*}
 F(8,3)
 &=\Phi_{11}(\zeta)
   (\zeta^6-\zeta^5+\zeta^4-\zeta^3+\zeta^2-2\zeta+3)=0,\\
 F(9,4)
 &=\Phi_{11}(\zeta)
   (\zeta^4-\zeta^3+\zeta^2+2\zeta-1)=0.
\end{align*}
Thus $|Z(F)|\ge6$.  Theorem~\ref{thm:five-point-six-zero} gives the reverse
inequality, so $|Z(F)|=6$.
\end{proof}



\begin{proposition}[Six-point threshold at ten zeros]
\label{prop:six-point-ten-threshold}
Let \(p\ge7\) and suppose \(|R|=6\).  If
\[
 |Z(F)|>10,
\]
then there exists a nonvertical affine line \(\ell\) such that
\begin{equation}
 F|_\ell\equiv0,
 \qquad
 \boxed{|Z(F)\setminus\ell|\le4.}
 \label{eq:six-point-ten-structure}
\end{equation}
\end{proposition}

\begin{proof}
Let \(r_Z\) be the zero-row rank.

If \(r_Z\le4\), a coordinate-hyperplane cut produces a nonzero
coefficient vector on at most five support points preserving every
zero.  Actual odd support at most five has at most six zeros by
Theorems~\ref{thm:three-point-exact} and
\ref{thm:five-point-six-zero}.  If the reduced actual support is four,
Proposition~\ref{prop:four-point-large-zero-structure} gives an affine
zero line and at most two preserved zeros off it.  Since more than
eight original zeros then lie on that line, the support-six
restriction cannot be nonzero there; Lemma~\ref{lem:affine-line-zero-bound}
forces the original sum to vanish identically on the same line.  The
support-two alternative is simpler.  Thus
\eqref{eq:six-point-ten-structure} holds in the rank-deficient branch.

Assume \(r_Z=5\), scale the coefficients into \(K=\Q(\omega)\), and
let \(n\le3\) be the first nonzero jet.

For \(n=0\) there are no zeros.  For \(n=1\), all exact zeros lie on
the affine line \(Q_0=0\); more than five force exact vanishing on a
nonvertical line.

For \(n=2\), if \(B\equiv0\) the vertical-fiber estimate gives at most
six zeros.  If \(B\not\equiv0\), factor
\[
 Q_0=D(a)C(a,b).
\]
A primitive nonlinear component has degree two, hence
\(\deg Q_0=n=2\); it gives at most six zeros by B\'ezout and
Lemma~\ref{lem:Q0-not-divide-Q1-universal}.  If \(C\) is a line, a
nonexact line contributes at most five zeros and there is at most one
exceptional vertical fiber, contributing at most three further zeros.
Hence more than ten zeros force an exact affine zero line; the off-line
vertical contribution is at most three.

Now let \(n=3\).  If \(B\equiv0\), the safe vertical estimate gives at
most nine zeros, except in the already-audited boundary configuration,
where the stronger estimate \(p-2\) applies.

Assume \(B\not\equiv0\).
Since \(n=3=\lfloor6/2\rfloor\),
Lemma~\ref{lem:maximal-order-full-degree} gives
\[
 \deg Q_0=3.
\]
Write
\[
 Q_0=D(a)C(a,b),
 \qquad
 m:=\deg D.
\]

If \(m=0\), then \(C=Q_0\) is a primitive cubic.
Lemma~\ref{lem:Q0-not-divide-Q1-universal} makes \(Q_0,Q_1\)
coprime.  If \(\deg Q_1\le3\), ordinary B\'ezout gives at most nine
affine common zeros.  If \(\deg Q_1=4\), then
\[
 \deg Q_0=3,\quad\deg_bQ_0=1,
 \qquad
 \deg Q_1=4,\quad\deg_bQ_1\le3,
\]
and Lemma~\ref{lem:bezout-infinity} gives at most
\[
 3\cdot4-(3-1)(4-3)=10
\]
affine common zeros.

If \(m=1\), then \(C\) is quadratic.  If \(C\nmid Q_1\), then
ordinary B\'ezout gives at most six exact zeros on \(C\) when
\(\deg Q_1\le3\); when \(\deg Q_1=4\),
Lemma~\ref{lem:bezout-infinity} gives at most
\(2\cdot4-(2-1)(4-3)=7\).  The exceptional vertical fiber contributes
at most three further zeros.

We claim \(C\mid Q_1\) is impossible.  Translate \(a\) so that
\(D(a)=a\), and write
\[
 C=A(a)+bB(a).
\]
Since \(\deg C=2\), one has \(\deg B\le1\).  If
\[
 Q_1=C(l_0+l_1b+l_2b^2),
\]
the recurrence gives
\begin{align*}
 6l_2B&=B+aB',\\
 2(A l_2+B l_1)&=A+aA'.
\end{align*}
The first equation forces \(B\) to be a monomial of degree \(0\) or
\(1\).  If \(\deg B=1\), primitivity gives \(A(0)\ne0\); evaluating the
second equation at \(a=0\) yields \(2l_2=1\), while the first gives
\(6l_2=2\), a contradiction in characteristic \(p\ge7\).
If \(B\) is constant, then \(l_2=1/6\), and
\[
 2Bl_1=aA'+\frac23A.
\]
Here \(\deg l_1\le1\), whereas \(\deg C=2\) forces \(\deg A=2\).  The
quadratic coefficient on the right is
\[
 \frac83[a^2]A\ne0,
\]
again a contradiction.  Thus the quadratic component cannot persist.

Finally let \(m=2\), so \(C\) is a nonvertical affine line.  If
\(C\nmid Q_1\), the line and \(Q_1\) have at most four common points,
while the multiplicity-weighted vertical budget is at most six; hence
\(|Z(F)|\le10\).

Assume \(C\mid Q_1\).  By the affine symmetries of
Lemma~\ref{lem:parabolic-affine-symmetry}, we may send \(C\) to the line
\(b=0\) without changing zero counts or vertical multiplicities.  Thus
\[
 Q_0=D(a)b,
 \qquad \deg D=2.
\]
The recurrence gives
\[
 Q_1=U(a)+V(a)b+\frac{D'(a)}6b^3.
\]
Since \(b\mid Q_1\), one has \(U\equiv0\), and therefore
\[
 Q_1=b\left(V(a)+\frac{D'(a)}6b^2\right).
\]
If \(D\) is irreducible over \(\F_p\), there is no exceptional vertical
fiber over \(\F_p\).  If \(D\) has two distinct roots, then at each root
\(\alpha\) the coefficient \(D'(\alpha)/6\) is nonzero, so the quotient is
a genuine quadratic in \(b\) and contributes at most two off-line zeros;
there are therefore at most four off-line zeros in total.

It remains to consider a double root \(\alpha\).  Translate \(a\) so that
\(\alpha=0\) and scale \(D\) so that \(D=a^2\).  If \(V(0)\ne0\), the
exceptional fiber contributes no off-line exact zero.  If \(V(0)=0\), then
\(Q_1(0,b)\equiv0\), so we use the next jet.  From
\(\partial_b^2Q_2=\partial_aQ_1\) one obtains
\[
 Q_2(0,b)=P_0+P_1b+\frac{V'(0)}6b^3+\frac1{60}b^5,
\]
whose leading coefficient is nonzero for every prime \(p\ge7\).  Thus the
exceptional fiber contains at most five exact zeros.

If \(F\) does not vanish identically on \(b=0\), the support-six line bound
gives at most five line zeros; together with the preceding fiber bound this
gives at most ten zeros.  Hence \(|Z(F)|>10\) forces \(F(a,0)=0\) for all
\(a\).  In that case \((0,0)\) is a root of the nonzero quintic
\(Q_2(0,b)\), so it has at most four further roots with \(b\ne0\).  Therefore
\[
 |Z(F)\setminus C|\le4.
\]
This completes the persistent-line case and the proof.
\end{proof}

\begin{theorem}[Seven-point refinement]
\label{thm:seven-point-ten}
Let \(p\ge11\) and suppose \(|R|=7\).  Then
\begin{equation}
 \boxed{|Z(F)|\le10.}
 \label{eq:seven-point-ten}
\end{equation}
\end{theorem}

\begin{proof}
Let \(r_Z\le6\) be the zero-row rank.

If \(r_Z\le4\), coordinate-hyperplane cuts reduce the common zero
system to actual support at most five.  Odd reductions have at most
six zeros, while a four-point reduction with a large zero set has an
affine zero line and at most two zeros off it.  Since the original
seven-point sum has at most six zeros on any affine line, these cases
give at most eight zeros.

If \(r_Z=5\), one coordinate cut gives a nonzero vector on at most six
support points preserving all zeros.  Supports at most five are
already controlled.  If the reduced actual support is six and more
than ten zeros are preserved,
Proposition~\ref{prop:six-point-ten-threshold} gives an affine zero
line with at most four zeros off it.  The original seven-point sum has
at most six zeros on that line, so
\[
 |Z(F)|\le6+4=10.
\]

It remains to consider \(r_Z=6\).  Scale into \(K\), and let \(n\le3\)
be the first nonzero jet.

For \(n=0,1\) the bound is immediate.  Suppose \(n=2\).  If
\(B\equiv0\), the vertical estimate gives at most six zeros.  Assume
\(B\not\equiv0\) and factor \(Q_0=D(a)C(a,b)\).  If \(C\) is a
primitive quadratic, then \(D\) is constant and
Lemma~\ref{lem:Q0-not-divide-Q1-universal} shows \(C\nmid Q_1\); B\'ezout
therefore gives at most six exact zeros.  If \(C\) is a line and
\(\deg D=0\), then \(Q_0\) itself is a nonvertical affine line, so every
exact zero lies on \(C\); odd-support line rigidity gives
\(|Z(F)|\le6\).  It remains to consider the line case \(\deg D=1\).
When \(C\nmid Q_1\), the line contains at most six exact zeros by
odd-support line rigidity, while the single exceptional vertical fiber
contains at most three, giving \(|Z(F)|\le9\).  When \(C\mid Q_1\), write
\(Q_1=CL\) with \(\deg L\le2\).  Let \(\alpha\in\F_p\) be the root of
the linear polynomial \(D\).  If
\(L(\alpha,b)\equiv0\), then \(Q_1(\alpha,b)\equiv0\), whereas the
recurrence gives
\[
 \partial_b^2Q_1(\alpha,b)
 =\partial_aQ_0(\alpha,b)
 =D'(\alpha)C(\alpha,b)\not\equiv0,
\]
a contradiction.  Thus \(L(\alpha,b)\) is nonzero and has at most two
roots, so there are at most two off-line zeros on the exceptional fiber and
\(|Z(F)|\le8\).  Hence in every \(n=2\) case
\(|Z(F)|\le9<10\).

Let \(n=3\).  If \(B\equiv0\), the vertical estimate gives at most nine
zeros.  For \(B\not\equiv0\), factor
\[
 Q_0=D(a)C(a,b).
\]
Since
\[
 n=3=\left\lfloor\frac72\right\rfloor,
\]
Lemma~\ref{lem:maximal-order-full-degree} gives
\[
 \deg Q_0=3.
\]

If \(\deg D=0\), the cubic \(Q_0\) and \(Q_1\) are coprime.  The same degree split used in the six-point proof (ordinary B\'ezout when \(\deg Q_1\le3\), and Lemma~\ref{lem:bezout-infinity} when \(\deg Q_1=4\)) gives at most ten affine zeros.

If \(\deg D=1\), the nonvertical component is quadratic.  If it does
not divide \(Q_1\), the same degree split as above gives at most seven
zeros on the component and the vertical fiber contributes at most three
further zeros.  The persistence alternative is impossible by
the quadratic calculation in the proof of
Proposition~\ref{prop:six-point-ten-threshold}.

If \(\deg D=2\), the nonvertical component is a line.  When it does
not divide \(Q_1\), the line contributes at most four zeros and the
vertical budget at most six, giving at most ten.

Assume now that the line divides \(Q_1\).  Normalize it to \(b=0\) as in
the persistent-line analysis of Proposition~\ref{prop:six-point-ten-threshold},
so \(Q_0=D(a)b\).  If \(D\) is irreducible or has two distinct roots over
\(\F_p\), the same explicit quotient formula
\[
 Q_1=b\left(V(a)+\frac{D'(a)}6b^2\right)
\]
gives at most four off-line zeros, while odd-support line rigidity gives at
most six line zeros; hence the total is at most ten.

Suppose finally that \(D\) has a double root.  Translate and scale so that
\(D=a^2\).  If the exceptional fiber is not annihilated by \(Q_1\), it
contributes no off-line zero.  Otherwise \(Q_1(0,b)\equiv0\), and the next
jet satisfies
\[
 Q_2(0,b)=P_0+P_1b+\frac{V'(0)}6b^3+\frac1{60}b^5,
\]
a nonzero polynomial of degree five.  Hence there are at most five exact
zeros on the exceptional vertical fiber.  If the line \(b=0\) contains at
most five exact zeros, the total is at most ten.  If it contains six exact
zeros, then the polynomial \(a\mapsto Q_2(a,0)\), whose degree is at most
\(5\), has six distinct roots and must vanish identically.  In particular
\(Q_2(0,0)=0\), so the nonzero quintic \(Q_2(0,b)\) has at most four roots
with \(b\ne0\).  Again the total is at most \(6+4=10\).

All cases give \eqref{eq:seven-point-ten}.
\end{proof}

\begin{remark}
The general odd-support theorem gives \(12\) for seven-point support.
Theorem~\ref{thm:seven-point-ten} shows that the uniform quadratic
bound is already nonsharp at seven-point support.  The present argument
lowers the bound from $12$ to $10$ but does not identify the exact uniform
maximum.  The improvement comes from forced intersection at infinity
together with the quotient control on persistent line components.
\end{remark}


\begin{remark}[Seven-point sharpness]
Theorem~\ref{thm:seven-point-ten} improves the general odd-support bound
from \(12\) to \(10\) at support seven.  We do not determine here whether
this constant is optimal for every prime \(p\ge11\); in particular, the
present argument does not determine the exact seven-point maximum when
\(p=11\).
\end{remark}

\subsection{Jet components and intersection theory}

\begin{theorem}[Full-degree jet profile and component erosion]
\label{thm:full-degree-component-erosion}
Let \(n\ge1\) be the first nonzero cyclotomic jet order, assume
\[
 \deg Q_0=n,
\]
and put
\[
 S_m:=\sum_{x\in R}c_{x,0}x^m.
\]
Then
\begin{equation}
 S_0=\cdots=S_{2n-2}=0,
 \label{eq:full-degree-moment-window}
\end{equation}
while
\[
 u:=S_{2n-1},\qquad v:=S_{2n}
\]
are not both zero.  For \(0\le j\le n-1\), the highest homogeneous
part of \(Q_j\) is
\begin{equation}
 Q_j^{[n+j]}(a,b)
 =
 \frac1{(n+j)!}
 \sum_{t=0}^{n+j}
 \binom{n+j}{t}
 a^tb^{n+j-t}S_{n+j+t}.
 \label{eq:full-degree-leading-form}
\end{equation}
Consequently,
\[
 \operatorname{ord}_aQ_j^{[n+j]}
 =
 \begin{cases}
 n-j-1,&u\ne0,\\
 n-j,&u=0.
 \end{cases}
\]
In particular, if \(C\) is an irreducible factor of \(Q_0\) of degree
\(d\ge2\), then
\begin{equation}
 \boxed{
 C\nmid Q_{\,n-d+1}.
 }
 \label{eq:full-degree-component-death}
\end{equation}
The indicated jet is always available.
\end{theorem}

\begin{proof}
For \(0\le q<n\), the highest homogeneous part of \(H_q\) is
\[
 \frac1{q!}
 \sum_{x\in R}c_{x,0}(ax^2+bx)^q.
\]
Since \(H_q=0\), the argument used in
Lemma~\ref{lem:first-nonzero-jet} gives
\[
 S_m=0
 \qquad(q\le m\le2q).
\]
The intervals \([q,2q]\), \(0\le q<n\), cover
\(0,\ldots,2n-2\), proving
\eqref{eq:full-degree-moment-window}.

Formula \eqref{eq:full-degree-leading-form} follows by taking the top
homogeneous part of \(H_{n+j}\).  For \(j=0\), the only potentially
nonzero contributions are \(S_{2n-1}\) and \(S_{2n}\).  If both
vanished, the homogeneous degree-\(n\) part of \(Q_0\) would vanish,
contrary to \(\deg Q_0=n\).  Thus \((u,v)\ne(0,0)\).

If \(u\ne0\), the first nonzero term in
\eqref{eq:full-degree-leading-form} occurs at
\[
 t=n-j-1
\]
and has \(b\)-exponent \(2j+1\).  If \(u=0\), then \(v\ne0\), and the
first nonzero term occurs at
\[
 t=n-j
\]
with \(b\)-exponent \(2j\).  All relevant binomial and factorial
coefficients are nonzero because
\[
 n+j\le2n-1\le p-2.
\]
This proves the exact \(a\)-adic orders.

Now let \(C\mid Q_0\) have degree \(d\ge2\).
The top form of \(Q_0\) is
\[
 Q_0^{[n]}
 =
 \frac{a^{n-1}}{n!}(nu\,b+va).
\]
If \(u\ne0\), every degree-\(d\) homogeneous divisor of this form is
divisible by at least \(a^{d-1}\); if \(u=0\), it is divisible by
\(a^d\).  Put
\[
 j=n-d+1.
\]
The preceding erosion law gives \(a\)-adic order \(d-2\) or \(d-1\),
respectively, for \(Q_j^{[n+j]}\).  Hence the highest homogeneous form
of \(C\) cannot divide that of \(Q_j\), and therefore \(C\nmid Q_j\).

Finally \(1\le j\le n-1\) and
\[
 n+j=2n-d+1\le2n-1\le p-2,
\]
so \(Q_j\) is available.
\end{proof}

\begin{theorem}[Full-degree odd-support bound]
\label{thm:full-degree-odd-bound}
Let
\[
 |R|=2r+1\le p-1,\qquad r\ge3,
\]
and suppose the exact zero-row system has rank \(2r\).
Scale the coefficients into \(K=\Q(\omega)\), let \(n\le r\) be the
first nonzero jet order, and assume
\[
 \deg Q_0=n.
\]
Then
\begin{equation}
 \boxed{|Z(F)|\le r(r+1).}
 \label{eq:full-degree-odd-bound}
\end{equation}
\end{theorem}

\begin{proof}
Write
\[
 Q_0=A(a)+bB(a).
\]
If \(B\equiv0\), all exact zeros lie on vertical fibers over roots of
\(A\).  Since \(2n\le2r\le p-3\), the safe fiber estimate applies to
every root and gives
\[
 |Z(F)|\le3n\le3r\le r(r+1).
\]

Assume \(B\not\equiv0\).  Put
\[
 D=\gcd(A,B),\qquad m=\deg D,
\]
and factor
\[
 Q_0=D(a)C(a,b),
\]
where \(C\) is primitive, irreducible, nonvertical, and
\[
 d:=\deg C=n-m.
\]
The exact zeros not on \(C\) lie on vertical fibers over roots of
\(D\).  If a root has multiplicity \(s\) in \(D\), then the common
order of \(A\) and \(B\) there is exactly \(s\), so the fiber contains
at most \(2s+1\) exact zeros.  Summing gives the vertical budget
\begin{equation}
 N_{\mathrm{vert}}\le3m.
 \label{eq:full-degree-vertical-budget}
\end{equation}

If \(d=1\), the curve \(C=0\) is an affine line and contains at most
\(2r\) exact zeros by the odd-support line bound.  Since
\(m=n-1\le r-1\),
\[
 |Z(F)|
 \le2r+3m
 \le5r-3
 \le r(r+1),
\]
where the last inequality holds for \(r\ge3\).

Assume \(d\ge2\).  By
Theorem~\ref{thm:full-degree-component-erosion},
\[
 C\nmid Q_{m+1},
\]
because
\[
 m+1=n-d+1.
\]
B\'ezout therefore gives
\[
 \#\bigl(Z(F)\cap C\bigr)
 \le
 d(n+m+1).
\]
Combining with \eqref{eq:full-degree-vertical-budget},
\[
 |Z(F)|
 \le
 (n-m)(n+m+1)+3m
 =
 n(n+1)+1-(m-1)^2.
 \label{eq:full-degree-almost}
\]
For \(m=0\) or \(m\ge2\), this is at most
\[
 n(n+1)\le r(r+1).
\]

It remains only \(m=1\).  If \(C\nmid Q_1\), use \(Q_1\) rather than
\(Q_2\):
\[
 \#\bigl(Z(F)\cap C\bigr)
 \le
 (n-1)(n+1)=n^2-1.
\]
Together with the vertical budget \(3\), this gives
\[
 |Z(F)|\le n^2+2\le n(n+1)\le r(r+1)
\]
because \(d=n-1\ge2\) implies \(n\ge3\).

If instead \(C\mid Q_1\), write
\[
 Q_1=C\,L.
\]
Since \(\deg C=n-1\) and \(\deg Q_1\le n+1\),
\[
 \deg L\le2.
\]
The component-erosion theorem gives
\[
 C\nmid Q_2.
\]
Hence
\[
 \#\bigl(Z(F)\cap C\bigr)
 \le
 (n-1)(n+2)
 =
 n^2+n-2.
\]
Let \(\alpha\in\F_p\) be the unique root of the linear polynomial
\(D\).  We claim that \(L(\alpha,b)\) is not the zero polynomial.  If
\(L(\alpha,b)\equiv0\), then \(Q_1(\alpha,b)\equiv0\), and hence
\(\partial_b^2Q_1(\alpha,b)=0\).  On the other hand, the jet recurrence gives
\[
 \partial_b^2Q_1(\alpha,b)
 =\partial_aQ_0(\alpha,b)
 =D'(\alpha)C(\alpha,b).
\]
Since \(D'(\alpha)\ne0\) and primitivity of \(C\) implies
\(C(\alpha,b)\not\equiv0\), this is a contradiction.  Thus every exact
zero on the exceptional vertical fiber but outside \(C\) is a root of the
nonzero polynomial \(L(\alpha,b)\), which has degree at most two.  There are
therefore at most two such points.  Consequently
\[
 |Z(F)|\le n(n+1)\le r(r+1).
\]
\end{proof}


\begin{lemma}[Persistence consumes vertical multiplicity]
\label{lem:persistence-vertical-multiplicity}
Let
\[
 Q_0=D(a)C(a,b),
 \qquad
 C=A(a)+bB(a),
\]
where \(C\) is primitive and nonvertical.  Suppose \(B\) is
nonconstant, \(\deg(DB)<p\), and
\[
 C\mid Q_1,\ldots,Q_J,
 \qquad
 2J+1<p.
\]
(The degree hypothesis is automatic in every application below, since
\(DB\) is the coefficient of \(b\) in \(Q_0\) and
\(\deg Q_0\le n\le(p-1)/2\).)
For \(j\ge1\), write
\[
 Q_j=C\,L_j
\]
and let \(u_j(a)\) be the coefficient of \(b^{2j}\) in \(L_j\);
put \(u_0=D\).  Then
\begin{equation}
 \boxed{
 B\,u_j
 =
 \frac{(DB)^{(j)}}{(2j+1)!}
 \qquad(0\le j\le J).
 }
 \label{eq:top-b-persistence}
\end{equation}
Consequently, for every root \(\beta\) of \(B\) in
\(\overline{\F}_p\),
\begin{equation}
 \operatorname{ord}_\beta D\ge J.
 \label{eq:vertical-multiplicity-consumption}
\end{equation}
In particular, if \(m=\deg D\), if the jets
\(Q_1,\ldots,Q_{m+1}\) are available, and if \(2m+3<p\), then
\begin{equation}
 \boxed{\exists\,j\in\{1,\ldots,m+1\}\ \text{such that}\ C\nmid Q_j.}
 \label{eq:nonconstant-B-death}
\end{equation}
Equivalently, a component satisfying these hypotheses must break by order
\(m+1\).
\end{lemma}

\begin{proof}
Since \(\deg_bQ_j\le2j+1\) and \(\deg_bC=1\),
\[
 \deg_bL_j\le2j.
\]
The coefficient of \(b^{2j-1}\) in
\[
 \partial_b^2(C L_j)=\partial_a(C L_{j-1})
\]
gives
\[
 (2j)(2j+1)B u_j=(B u_{j-1})'.
\]
Induction yields \eqref{eq:top-b-persistence}; the factorials are
invertible under \(2J+1<p\).

Let \(\beta\) be a root of \(B\), with
\[
 e=\operatorname{ord}_\beta B,\qquad
 s=\operatorname{ord}_\beta D.
\]
We prove \(s\ge j\) inductively for \(1\le j\le J\).

For \(j=1\), the polynomial \(DB\) has local order \(s+e<p\).
Its derivative therefore has exact local order \(s+e-1\), because the
leading falling-factorial coefficient is nonzero in \(\F_p\).
The divisibility
\[
 B\mid(DB)'
\]
forces
\[
 s+e-1\ge e,
\]
hence \(s\ge1\).

Assume now \(s\ge j-1\).  Since \(e\ge1\),
\[
 j\le s+e.
\]
Thus the \(j\)-th derivative of \(DB\) has exact local order
\[
 s+e-j;
\]
again the leading falling factorial is nonzero because
\(\deg(DB)<p\).  The divisibility
\[
 B\mid(DB)^{(j)}
\]
then gives
\[
 s+e-j\ge e,
\]
so \(s\ge j\).  This proves
\eqref{eq:vertical-multiplicity-consumption}.  For the final assertion,
assume for contradiction that \(C\mid Q_j\) for every
\(1\le j\le m+1\).  Applying the preceding conclusion with
\(J=m+1\) to any root \(\beta\in\overline{\F}_p\) of the nonconstant
polynomial \(B\) gives
\[
 \operatorname{ord}_\beta D\ge m+1,
\]
contrary to \(\deg D=m\).  Hence persistence through all of
\(Q_1,\ldots,Q_{m+1}\) is impossible, which proves
\eqref{eq:nonconstant-B-death}.
\end{proof}

\begin{theorem}[Degree-drop branch with nonconstant graph coefficient]
\label{thm:degree-drop-nonconstant-B}
Let
\[
 |R|=2r+1\le p-1,\qquad r\ge3,
\]
and suppose the exact zero-row system has rank \(2r\).
Scale the coefficients into \(K=\Q(\omega)\), let \(n\le r\) be the
first nonzero jet order, and assume
\[
 q:=\deg Q_0<n.
\]
Factor
\[
 Q_0=D(a)C(a,b),
 \qquad
 C=A(a)+bB(a),
\]
with \(C\) primitive and nonvertical.  If \(B\) is nonconstant, then
\begin{equation}
 \boxed{|Z(F)|\le r(r+1).}
 \label{eq:degree-drop-nonconstant-B}
\end{equation}
\end{theorem}

\begin{proof}
Put
\[
 m=\deg D,\qquad d=\deg C,
\]
so
\[
 m+d=q\le n-1.
\]
Since \(B\) is nonconstant, \(d\ge2\), hence
\[
 m\le n-3.
\]
The vertical exceptional fibers contribute at most \(3m\) exact
zeros.

Because \(m\le n-3\),
\[
 n+m+1\le2n-2\le p-3
 \qquad\text{and}\qquad
 2(m+1)+1\le2n-3<p.
\]
Hence \(Q_{m+1}\) is available and the factorial condition in
Lemma~\ref{lem:persistence-vertical-multiplicity} is satisfied through order
\(m+1\).  If \(C\) divided every one of
\(Q_1,\ldots,Q_{m+1}\), that lemma would force a root
\(\beta\in\overline{\F}_p\) of the nonconstant polynomial \(B\) to
satisfy
\[
 \operatorname{ord}_\beta D\ge m+1,
\]
which is impossible because \(\deg D=m\).  Thus there is a first index
\(J\ge1\) for which
\[
 C\nmid Q_J,
\]
and necessarily
\[
 J\le m+1.
\]
In particular,
\[
 n+J\le n+m+1\le2n-2\le p-3,
\]
so every jet used below is available.

Every exact zero on \(C\) also annihilates \(Q_J\).
Since \(C\) is irreducible and does not divide \(Q_J\), B\'ezout gives
\[
 \#\bigl(Z(F)\cap C\bigr)
 \le
 d(n+J)
 \le
 d(n+m+1).
\]
Using
\[
 d\le n-m-1,
\]
we obtain
\[
 \begin{aligned}
 |Z(F)|
 &\le
 (n-m-1)(n+m+1)+3m\\
 &=n^2-m^2+m-1\\
 &\le n^2
 \le r(r+1).
 \end{aligned}
\]
\end{proof}




\begin{lemma}[Degree-drop polynomial graphs break by the second jet]
\label{lem:degree-drop-graph-break}
Let \(n\) be the first nonzero jet order and suppose
\[
 Q_0=D(a)\bigl(A(a)+b\bigr),
\]
where
\[
 m:=\deg D,\qquad d:=\deg A\ge2,\qquad
 m+d<n.
\]
Put
\[
 \delta:=n-m-d\ge1.
\]
If
\[
 A(a)+b\mid Q_1
 \qquad\text{and}\qquad
 d\ge\delta+2,
\]
then
\begin{equation}
 \boxed{A(a)+b\nmid Q_2.}
 \label{eq:degree-drop-graph-break}
\end{equation}
\end{lemma}

\begin{proof}
Write
\[
 C=A+b,\qquad Q_1=C\,L,
 \qquad
 L=l_0+l_1b+l_2b^2.
\]
Since
\[
 \deg L\le n+1-d=m+\delta+1,
\]
one has
\[
 \deg l_1\le m+\delta.
\]
The recurrence
\[
 \partial_b^2Q_1=\partial_aQ_0
\]
gives
\[
 l_2=\frac{D'}6,
 \qquad
 l_1=\frac12DA'+\frac13D'A.
\]
If \(D\) and \(A\) have leading coefficients \(D_m,A_d\), then the
coefficient of \(a^{m+d-1}\) in \(l_1\) is
\[
 \frac{D_mA_d}{6}(3d+2m).
\]
Because \(d\ge\delta+2\), the nominal degree \(m+d-1\) is strictly
larger than the allowed degree \(m+\delta\).  Hence persistence through
\(Q_1\) forces the resonance
\begin{equation}
 3d+2m=0
 \qquad\text{in }\F_p.
 \label{eq:first-graph-resonance}
\end{equation}

Suppose, for contradiction, that \(C\mid Q_2\).  Write
\[
 Q_2=C\,M,
 \qquad
 M=m_0+m_1b+m_2b^2+m_3b^3+m_4b^4.
\]
The hypotheses force \(n\ge4\).  In the standing parabolic range
\(n\le\lfloor |R|/2\rfloor\le(p-1)/2\), this implies \(p\ge11\) and
\(n+2\le p-2\).  Thus \(Q_2\) is available, and the denominators
\(6\) and \(120\) appearing below are invertible in \(\F_p\).
Comparing the \(b^3\)- and \(b^2\)-coefficients in
\[
 \partial_b^2Q_2=\partial_aQ_1
\]
gives
\[
 m_4=\frac{D''}{120},
\]
and
\[
 m_3=
 \frac{
 4AD''+10D'A'+5DA''
 }{120}.
\]
Since
\[
 \deg M\le n+2-d=m+\delta+2,
\]
one has
\[
 \deg m_3\le m+\delta-1.
\]
But the nominal degree of the displayed numerator is \(m+d-2\),
strictly larger because \(d\ge\delta+2\).  Therefore its leading
coefficient must vanish:
\begin{equation}
 4m(m-1)+10md+5d(d-1)=0
 \qquad\text{in }\F_p.
 \label{eq:second-graph-resonance}
\end{equation}
Using \eqref{eq:first-graph-resonance} to substitute
\(m=-3d/2\), the left-hand side of
\eqref{eq:second-graph-resonance} becomes
\[
 d(1-d).
\]
Here \(d\ge3\) and \(d<p\), so this is nonzero in \(\F_p\), a
contradiction.
\end{proof}

\subsection{Restriction to rational components}

The early-jet regime naturally leads to restrictions of $F$ to curves that are linear in $b$.

\begin{theorem}[Rational-component restriction]
\label{thm:rational-component-restriction}
Let
\[
 C(a,b)=A(a)+bB(a)\in\F_p[a,b]
\]
be irreducible and nonvertical of degree $d\ge2$.  For actual support $|R|=k\le p-1$,
\begin{equation}
 \boxed{\#\bigl(Z(F)\cap C(\F_p)\bigr)\le d(k-1).}
 \label{eq:rational-component-restriction}
\end{equation}
In particular, an irreducible conic contains at most $2(k-1)$ exact zeros.
\end{theorem}

\begin{proof}
We argue by induction on $k$.  For $k=1$ the sum has no zeros, so the claim is immediate.  Assume $k\ge2$, put
\[
 Z_C:=Z(F)\cap C(\F_p),
\]
and let $r_C$ be the rank of the evaluation rows indexed by $Z_C$.
Since the coefficient vector lies in their common kernel, $r_C\le k-1$.
If $r_C\le k-2$, Lemma~\ref{lem:complex-to-cyclotomic} gives a nonzero
coefficient vector on a strictly smaller actual support preserving every
zero in $Z_C$; the induction hypothesis applies.  If $r_C=k-1$, the same
lemma scales the coefficient vector into $K$.  Thus it remains only to
treat primitive cyclotomic coefficients.

If $B$ is constant, write the curve as $b=P(a)$ with $\deg P=d$.  For
\[
 \mathcal Y(T;a)=\sum_{x\in R}C_x(T)e^{T(ax^2+P(a)x)}=\sum J_j(a)T^j,
\]
one has $\deg J_j\le dj$.  If $J_0=\cdots=J_{k-1}=0$, the coefficient of $a^{dj}$ in $J_j$ gives $\sum_xc_{x,0}x^j=0$ for $0\le j\le k-1$, contradicting Vandermonde.  Thus the first nonzero $J_n$ has $n\le k-1$, and all exact zeros are roots of a nonzero polynomial of degree at most $dn\le d(k-1)$.

Assume $B$ is nonconstant.  Since $\gcd(A,B)=1$, affine points satisfy $B(a)\ne0$ and
\[
 b=-A(a)/B(a).
\]
Set $H_x(a)=ax^2-A(a)x/B(a)$ and define $\mathcal Y$ as above.  Then $B(a)^jJ_j(a)$ has degree at most $dj$.  Choose a root $\alpha$ of $B$ of multiplicity $m$ in $\overline{\F}_p$; because $A(\alpha)\ne0$, the principal part of $H_x$ at $\alpha$ is $\gamma x(a-\alpha)^{-m}$.  If $J_0,\ldots,J_{k-1}$ vanished, then $J_0=0$ gives $\sum_xc_{x,0}=0$, while for $1\le j\le k-1$ the highest-pole coefficient of $J_j$ gives $\sum_xc_{x,0}x^j=0$.  These $k$ moments contradict Vandermonde invertibility.  Therefore $n\le k-1$, and every exact zero on $C(\F_p)$ is a root of the nonzero polynomial $B^nJ_n$ of degree at most $dn$.
\end{proof}

\begin{remark}[A graph-spectrum statement]
\label{rem:graph-spectrum-rational}
The proof of Theorem~\ref{thm:rational-component-restriction} does not
use the quadratic term \(x^2\) in any essential way.  Replacing
\(ax^2\) by \(a\phi(x)\) for an arbitrary function
\(\phi:\F_p\to\F_p\) leaves the pole argument and the Vandermonde
moments unchanged.  Consequently the same estimate
\[
 \#\bigl(Z(F_\phi)\cap C(\F_p)\bigr)\le d(k-1)
\]
holds for primitive graph-spectrum data
\[
 F_\phi(a,b)=\sum_{x\in R}c_x\omega^{a\phi(x)+bx}.
\]
This is the part of the component method that is genuinely
graph-theoretic rather than quadratic.
\end{remark}

\begin{proposition}[Degree-drop graph estimate]
\label{prop:degree-drop-graph-estimate}
Let the actual support have size \(k\le2t+1\), with \(t\ge3\),
and suppose the exact zero-row system has rank \(k-1\).  After scaling
the coefficient vector into \(K=\Q(\omega)\) as in
Lemma~\ref{lem:complex-to-cyclotomic}, let \(n\le\lfloor k/2\rfloor\le t\)
be the first nonzero cyclotomic jet order and suppose
\[
 Q_0=D(a)\bigl(A(a)+b\bigr),
 \qquad
 m:=\deg D,\quad d:=\deg A\ge2,\quad
 m+d<n.
\]
Then
\begin{equation}
 \boxed{|Z(F)|\le t(t+1).}
 \label{eq:degree-drop-graph-estimate}
\end{equation}
\end{proposition}

\begin{proof}
Put
\[
 \delta=n-m-d\ge1,
 \qquad
 C=A+b.
\]
The vertical exceptional fibers contribute at most \(3m\) exact zeros.

If \(C\nmid Q_1\), B\'ezout gives
\[
 \#(Z(F)\cap C)\le d(n+1).
\]
Since
\[
 d\le n-m-1,
\]
\[
 |Z(F)|
 \le
 (n-m-1)(n+1)+3m
 =
 n^2-1+m(2-n)
 \le n^2
 \le t(t+1).
\]

Assume \(C\mid Q_1\).
First suppose
\[
 d\le\delta+1.
\]
The rational-component restriction theorem gives
\[
 \#(Z(F)\cap C)
 \le
 d(|R|-1)
 \le2td.
\]
Also
\[
 n=m+d+\delta\le t,
 \qquad
 \delta\ge d-1,
\]
so
\[
 m+2d\le t+1.
\]
Therefore
\[
 |Z(F)|
 \le
 2td+3m
 \le
 2td+3(t+1-2d).
\]
The difference from \(t(t+1)\) is
\[
 (t-3)(t+1-2d)\ge0.
\]
Thus \eqref{eq:degree-drop-graph-estimate} holds.

It remains to consider
\[
 d\ge\delta+2.
\]
Lemma~\ref{lem:degree-drop-graph-break} gives \(C\nmid Q_2\).
Hence
\[
 \#(Z(F)\cap C)\le d(n+2).
\]
For the vertical exceptional fibers we deliberately use the robust
multiplicity budget from the main zero-bound proof.  If \(\alpha\) is a
root of \(D\) of multiplicity \(s\), the common multiplicity of the two
coefficients of the leading pencil at \(\alpha\) is exactly \(s\); hence
that fiber contains at most \(2s+1\) exact zeros.  Summing over the roots of
\(D\) gives at most \(3m\) vertical zeros.  This estimate remains valid even
when \(D\) has repeated roots, where the quotient \(L(\alpha,b)\) may
degenerate.
Using again \(d\le n-m-1\),
\[
 \begin{aligned}
 |Z(F)|
 &\le d(n+2)+3m\\
 &\le(n-m-1)(n+2)+3m\\
 &=n(n+1)-2-m(n-1)\\
 &\le n(n+1)\\
 &\le t(t+1).
 \end{aligned}
\]

\end{proof}

\begin{proposition}[Cyclotomic nonlinear-component estimates]
\label{prop:cyclotomic-nonlinear-estimates}
Let \(n\) be the first nonzero cyclotomic jet and factor
\[
 Q_0=D(a)C(a,b),
\]
where \(C\) is primitive, nonvertical, irreducible, and
\[
 d:=\deg C\ge2.
\]
Then the following hold.

\begin{enumerate}
\item If \(\deg Q_0=n\), then
\[
 |Z(F)|\le n(n+1).
\]

\item If \(\deg Q_0<n\) and the coefficient of \(b\) in \(C\) is
nonconstant, then
\[
 |Z(F)|\le n^2.
\]
\end{enumerate}
\end{proposition}

\begin{proof}
For the first assertion put \(m=\deg D\), so \(d=n-m\).
The full-degree component-erosion theorem gives
\[
 C\nmid Q_{m+1}.
\]
The vertical fibers contribute at most \(3m\), while B\'ezout gives
\[
 \#(Z(F)\cap C)\le(n-m)(n+m+1).
\]
Thus
\[
 |Z(F)|
 \le
 n(n+1)+1-(m-1)^2.
\]
Only \(m=1\) requires a one-point saving.
If \(C\nmid Q_1\), use \(Q_1\) to get
\[
 (n-1)(n+1)+3=n^2+2\le n(n+1).
\]
If \(C\mid Q_1\), then \(C\nmid Q_2\), and write \(Q_1=CL\) with
\(\deg L\le2\).  Here \(m=1\), so \(D\) has a unique simple root
\(\alpha\in\F_p\).  We claim that \(L(\alpha,b)\) is not the zero
polynomial.  Indeed, if it were, then \(Q_1(\alpha,b)\equiv0\) and hence
\(\partial_b^2Q_1(\alpha,b)\equiv0\).  But the recurrence gives
\[
 \partial_b^2Q_1(\alpha,b)
 =\partial_a(DC)(\alpha,b)
 =D'(\alpha)C(\alpha,b).
\]
The first factor is nonzero because the root is simple, and the second is a
nonzero polynomial in \(b\) because \(C\) is primitive.  This contradiction
proves the claim.  Consequently \(L(\alpha,b)\), of degree at most two in
\(b\), has at most two roots.  Thus at most two exact zeros on the exceptional
vertical fiber lie outside \(C\), and
\[
 (n-1)(n+2)+2=n(n+1).
\]

For the second assertion let
\[
 q=\deg Q_0<n,\qquad m=\deg D,\qquad d=q-m.
\]
Because the coefficient of \(b\) in \(C\) is nonconstant, \(d\ge2\),
and the degree drop gives \(m\le n-3\).  Hence
\[
 n+m+1\le2n-2\le p-3,
 \qquad
 2(m+1)+1\le2n-3<p.
\]
Moreover the coefficient of \(b\) in \(Q_0\) has degree below \(p\).
Thus every hypothesis of
Lemma~\ref{lem:persistence-vertical-multiplicity} is satisfied through
\(Q_{m+1}\), and that lemma gives a first failure index
\(J\le m+1\) with \(C\nmid Q_J\).  Using
\(n+J\le n+m+1\) in the subsequent B\'ezout estimate, and since
\[
 d\le n-m-1,
\]
B\'ezout plus the vertical budget gives
\[
 \begin{aligned}
 |Z(F)|
 &\le
 (n-m-1)(n+m+1)+3m\\
 &=n^2-m^2+m-1\\
 &\le n^2.
 \end{aligned}
\]
\end{proof}

\begin{theorem}[Parity-sensitive zero theorem]
\label{thm:parity-sensitive-zero}
Let \(p\ge5\) be prime.

\begin{enumerate}
\item If the actual support has odd size
\[
 |R|=2r+1\le p-1,
\]
then
\begin{equation}
 \boxed{|Z(F)|\le r(r+1).}
 \label{eq:odd-uniform-zero}
\end{equation}

\item If the actual support has even size
\[
 |R|=2r\le p-1
\]
and
\begin{equation}
 |Z(F)|>r(r+1),
 \label{eq:even-large-threshold}
\end{equation}
then there exists a nonvertical affine line \(\ell\) such that
\begin{equation}
 F|_\ell\equiv0,
 \qquad
 \boxed{|Z(F)\setminus\ell|\le r(r-1).}
 \label{eq:even-line-structure}
\end{equation}
\end{enumerate}
\end{theorem}

\begin{proof}
We prove the two assertions simultaneously by induction on \(r\).
For \(r=0\), the odd assertion is the one-point case, and the sum has no
zeros.  For \(r=1\), the odd assertion is
Theorem~\ref{thm:three-point-exact}; a two-point sum has either no
zeros or an affine zero line.  For \(r=2\), the odd assertion is
Theorem~\ref{thm:five-point-six-zero}, and the even assertion is
Proposition~\ref{prop:four-point-large-zero-structure}.

Assume \(r\ge3\), and first prove the even assertion for support
\(2r\).

Let \(r_Z\) be the zero-row rank.  If \(r_Z\le2r-2\), a
coordinate-hyperplane cut gives a nonzero coefficient vector on at
most \(2r-1\) support points preserving every zero.

If the reduced actual support is one, the reduced sum has no zeros,
which is impossible because it preserves \(Z(F)\).  If the reduced actual
support is odd and at least three, say \(2s+1\) with \(1\le s\le r-1\),
the induction hypothesis gives
\[
 |Z(F)|\le s(s+1)\le r(r-1),
\]
contrary to \eqref{eq:even-large-threshold}.

If the reduced support is even, say \(2s\) with \(s\le r-1\), the
induction hypothesis yields an affine zero line for the reduced sum
and at most \(s(s-1)\) preserved zeros off it.  Hence more than
\[
 r(r+1)-s(s-1)
 \ge
 r(r+1)-(r-1)(r-2)
 =
4r-2
\]
zeros of the original sum lie on that line.  This exceeds \(2r-1\),
so Lemma~\ref{lem:affine-line-zero-bound} forces the original
support-\(2r\) sum to vanish identically on the same line.
Furthermore
\[
 |Z(F)\setminus\ell|
 \le
 s(s-1)
 \le
 r(r-1).
\]

It remains to treat the full-rank even branch \(r_Z=2r-1\).
Scale the coefficients into \(K=\Q(\omega)\), and let \(n\le r\) be
the first nonzero jet.

If \(B\equiv0\), the vertical-fiber estimate in the proof of the main
zero theorem gives at most \(3n\) zeros, except for its single boundary
configuration, where the stronger bound \(p-2\) applies.  In either
case the total is at most \(r(r+1)\).

Assume \(B\not\equiv0\), and factor
\[
 Q_0=D(a)C(a,b),
\]
where \(C\) is primitive and nonvertical.

If \(\deg C=1\), then \(C=0\) is an affine line.  If \(F\) does not
vanish identically there, the line contains at most \(2r-1\) exact
zeros.  Since \(\deg D\le n-1\le r-1\), the vertical budget is at most
\(3(r-1)\), and
\[
 (2r-1)+3(r-1)=5r-4\le r(r+1).
\]
Thus a zero set satisfying \eqref{eq:even-large-threshold} forces
\(F|_C\equiv0\).  All zeros off \(C\) then lie on vertical exceptional
fibers, so
\[
 |Z(F)\setminus C|
 \le3(r-1)
 \le r(r-1).
\]

Suppose \(\deg C\ge2\).  If \(\deg Q_0=n\), apply
Proposition~\ref{prop:cyclotomic-nonlinear-estimates}(1).
If \(\deg Q_0<n\) and the coefficient of \(b\) in \(C\) is
nonconstant, apply part (2).  If that coefficient is constant, apply
Proposition~\ref{prop:degree-drop-graph-estimate}.  Every case gives
\[
 |Z(F)|\le r(r+1),
\]
contrary to \eqref{eq:even-large-threshold}.  This proves the even
assertion.

Now prove the odd assertion for support \(2r+1\).
If \(r_Z\le2r-1\), a coordinate cut gives a vector on at most \(2r\)
support points preserving all zeros.

A reduced support of size one has no zeros.  A reduced odd support
\(2s+1\) with \(1\le s\le r-1\) has at most
\(s(s+1)<r(r+1)\) zeros by induction.
For a reduced even support \(2s\), \(s\le r\), the just-proved even
assertion gives an affine zero line and at most
\[
 s(s-1)\le r(r-1)
\]
zeros off it whenever more than \(r(r+1)\) zeros are preserved.
But the original odd-support sum has at most \(2r\) zeros on any
affine line.  Hence
\[
 |Z(F)|
 \le
 2r+r(r-1)
 =
 r(r+1),
\]
a contradiction.

Finally assume the odd zero-row system has full rank \(2r\).
Again scale into \(K\) and let \(n\le r\) be the first nonzero jet.
If \(B\equiv0\), the vertical estimate gives
\[
 |Z(F)|\le3n\le r(r+1).
\]
For \(B\not\equiv0\), factor \(Q_0=D C\).
If \(\deg C=1\), the odd line bound gives at most \(2r\) zeros on
\(C\), while the vertical budget is at most \(3(r-1)\); thus
\[
 |Z(F)|\le5r-3\le r(r+1).
\]
If \(\deg C\ge2\), there are three subcases, exactly as in the even
full-rank branch.  If \(\deg Q_0=n\),
Proposition~\ref{prop:cyclotomic-nonlinear-estimates}(1) gives
\[
 |Z(F)|\le n(n+1)\le r(r+1).
\]
If \(\deg Q_0<n\) and the coefficient of \(b\) in \(C\) is
nonconstant, Proposition~\ref{prop:cyclotomic-nonlinear-estimates}(2) gives
\[
 |Z(F)|\le n^2\le r(r+1).
\]
Finally, if that coefficient is constant, primitive normalization writes
\(C=A(a)+b\), and Proposition~\ref{prop:degree-drop-graph-estimate},
with \(t=r\), gives directly
\[
 |Z(F)|\le r(r+1).
\]
Thus every full-rank case is covered, and the induction is complete.
\end{proof}

\begin{corollary}[Combined odd-support bound]
\label{cor:combined-odd-support}
If
\[
 |R|=2r+1\le p-1,
\]
then
\begin{equation}
 \boxed{
 |Z(F)|
 \le
 \min\bigl\{
 p+2r-2,\,
 r(r+1)
 \bigr\}.
 }
 \label{eq:combined-odd-support}
\end{equation}
\end{corollary}

\begin{proof}
Combine Theorem~\ref{thm:parity-sensitive-zero} with the odd case of
Theorem~\ref{thm:parabolic-zero-rigorous}.
\end{proof}

\begin{corollary}[Maximal-jet branch]
\label{cor:maximal-jet-odd-bound}
Let
\[
 |R|=2r+1\le p-1,\qquad r\ge3,
\]
and suppose the exact zero-row system has rank \(2r\).  Scale the
coefficients into \(K=\Q(\omega)\), and let \(n\le r\) be the first
nonzero jet order.  If \(n=r\), then
\[
 |Z(F)|\le r(r+1).
\]
\end{corollary}

\begin{proof}
The vanishing \(H_0,\ldots,H_{r-1}=0\) gives
\[
 S_0=\cdots=S_{2r-2}=0.
\]
If \(\deg Q_0<r\), then also
\[
 S_{2r-1}=S_{2r}=0.
\]
The \((2r+1)\times(2r+1)\) Vandermonde moment matrix on the distinct
support points would then force every primitive residue
\(c_{x,0}\) to vanish, a contradiction.  Hence \(\deg Q_0=r\), and
Theorem~\ref{thm:full-degree-odd-bound} applies.
\end{proof}

\begin{corollary}[Uniform odd-support rigidity]
\label{cor:odd-uniform-summary}
For every odd actual support
\[
 |R|=2r+1\le p-1,
\]
one has
\[
 |Z(F)|\le r(r+1).
\]
For \(r=1\) the constant \(2\) is attained by Theorem~\ref{thm:three-point-exact}, and for \(r=2\) the constant \(6\) is attained by Proposition~\ref{prop:five-point-sharpness}.
\end{corollary}

\section{Sharp uncertainty consequences}
\label{sec:consequences}

\begin{corollary}[Restricted Fourier uncertainty on the parabola]
\label{cor:restricted-parabola-rigorous}
Let the Fourier transform on \(\F_p^2\) be
\[
 \widehat f(\xi,\eta)
 =
 \sum_{a,b\in\F_p}
 f(a,b)\omega^{-(a\xi+b\eta)}.
\]
Suppose \(f\ne0\) and
\[
 \supp\widehat f
 \subseteq
 \Gamma:=\{(x^2,x):x\in\F_p\},
 \qquad
 |\supp\widehat f|=k\le p-1.
\]
If \(k=2r\), then
\[
 |\supp f|
 \ge
 p^2-p-2r+2.
\]
If \(k=2r+1\), then
\begin{equation}
 \boxed{
 |\supp f|
 \ge
 \max\{
 p^2-p-2r+2,\,
 p^2-r(r+1)
 \}.
 }
 \label{eq:restricted-parabola-odd-improved}
\end{equation}
In particular, for fixed odd \(k\), the complement of \(\supp f\) has
cardinality bounded independently of \(p\).
\end{corollary}

\begin{proof}
Write
\[
 \widehat f(x^2,x)=d_x
 \qquad(x\in R),
\]
with \(d_x\ne0\).  Fourier inversion gives
\[
 f(a,b)
 =
 \frac1{p^2}
 \sum_{x\in R}d_x\omega^{ax^2+bx}.
\]
Thus the zeros of \(f\) are exactly the zeros of the corresponding
parabolic sum.  Apply Theorem~\ref{thm:parabolic-zero-rigorous} in the
even case and Corollary~\ref{cor:combined-odd-support} in the odd case,
then subtract the zero count from \(p^2\).
\end{proof}

\begin{theorem}[Exact complete-MUB uncertainty]
\label{thm:exact-mub-rigorous}
For every prime \(p\ge5\) and every nonzero
\(\psi\in\C^p\),
\begin{equation}
 \boxed{
 \mathcal S_p(\psi)\ge p^2-p+2.
 }
 \label{eq:exact-mub-bound-rigorous}
\end{equation}
The bound is attained.  Hence the exact minimum total support is
\begin{equation}
 \boxed{
 T_s(p)=p^2-p+2.
 }
 \label{eq:Ts-rigorous}
\end{equation}
\end{theorem}

\begin{proof}
Among the \(p+1\) standard MUBs, choose one in which \(\psi\) has
minimum support, and denote that support size by \(k\).
By Lemma~\ref{lem:standard-mub-transitivity}, after applying a unitary
which merely permutes the standard MUBs, we may assume that this basis is
the computational basis.  Thus
\[
 R:=\supp_{\mathcal B_\infty}(\psi)
\]
has size \(k\).

If \(k=p\), minimality implies that every one of the \(p+1\) support
sizes equals \(p\), so
\[
 \mathcal S_p(\psi)=p(p+1)>p^2-p+2.
\]

If \(k=1\), then \(\psi\) is a scalar multiple of one computational
basis vector.  Mutual unbiasedness implies that all \(p\) coordinates
in each of the other \(p\) bases are nonzero.  Hence
\[
 \mathcal S_p(\psi)=1+p^2>p^2-p+2.
\]

Assume now
\[
 2\le k\le p-1.
\]
Write
\[
 \psi=\sum_{x\in R}c_x e_x,
 \qquad c_x\ne0.
\]
For \(a,b\in\F_p\),
\[
 \langle\phi_{a,b},\psi\rangle
 =
 p^{-1/2}
 \sum_{x\in R}
 c_x\omega^{-a x^2-b x}.
\]
Replacing \((a,b)\) by \((-a,-b)\), which is a bijection of
\(\F_p^2\), shows that the total number \(N\) of zero outcomes among the
\(p\) finite quadratic bases is
\[
 N
 =
 \#\left\{
 (a,b)\in\F_p^2:
 \sum_{x\in R}c_x\omega^{a x^2+b x}=0
 \right\}.
\]
By Theorem~\ref{thm:parabolic-zero-rigorous},
\[
 N
 \le
 p+2\left\lfloor\frac{k}{2}\right\rfloor-2.
\]
The computational basis contributes support \(k\), while the \(p\)
finite bases contain \(p^2-N\) nonzero outcomes.  Therefore
\[
 \begin{aligned}
 \mathcal S_p(\psi)
 &=k+p^2-N\\
 &\ge
 k+p^2-
 \left(
 p+2\left\lfloor\frac{k}{2}\right\rfloor-2
 \right)\\
 &=
 p^2-p+
 k-2\left\lfloor\frac{k}{2}\right\rfloor+2\\
 &\ge p^2-p+2.
 \end{aligned}
\]
For odd \(k\), the displayed estimate is in fact
\(p^2-p+3\).

It remains to show that \eqref{eq:exact-mub-bound-rigorous} is attained.
Choose distinct \(r,s\in\F_p\) and \(t\in\F_p\), and put
\[
 \psi=e_r-\omega^t e_s.
\]
Its computational support is \(2\).  In the basis \(\mathcal B_a\),
\[
 \langle\phi_{a,b},\psi\rangle
 =
 p^{-1/2}
 \left(
 \omega^{-ar^2-br}
 -
 \omega^t\omega^{-as^2-bs}
 \right).
\]
This coefficient vanishes if and only if
\[
 -ar^2-br
 \equiv
 t-as^2-bs
 \pmod p,
\]
equivalently,
\begin{equation}
 b(r-s)
 =
 -a(r^2-s^2)-t.
 \label{eq:unique-zero-each-basis}
\end{equation}
Since \(r-s\ne0\), for every fixed \(a\in\F_p\) equation
\eqref{eq:unique-zero-each-basis} has exactly one solution
\(b\in\F_p\).  Thus each of the \(p\) finite quadratic bases has exactly
one zero coefficient and therefore support \(p-1\).  Hence
\[
 \mathcal S_p(\psi)
 =
 2+p(p-1)
 =
 p^2-p+2.
\]
This proves both sharpness and \eqref{eq:Ts-rigorous}.
\end{proof}

\begin{corollary}[Parity gap]
\label{cor:parity-gap-rigorous}
If a state has minimum support \(k\ge3\) among the standard MUBs and
\(k\) is odd, then
\[
 \mathcal S_p(\psi)\ge p^2-p+3.
\]
In particular, every state attaining
\(T_s(p)=p^2-p+2\) has even minimum support.
\end{corollary}


\section{Rigidity of equality}
\label{sec:extremizer-structure}

To avoid a sign ambiguity in the equality theory, we distinguish the two
coordinate systems from this point onward.  Lowercase \((a,b)\) always denotes
the coordinates of the parabolic sum \(F(a,b)\).  When referring to the
physical MUB measurement labels, we write \((A,\beta)\), where
\(\phi_{A,\beta}\in\mathcal B_A\).  By
\eqref{eq:mub-parabola-label-map},
\begin{equation}
 A=-a,\qquad \beta=-b.
 \label{eq:F-to-physical-labels}
\end{equation}

The sharp constant in Theorem~\ref{thm:exact-mub-rigorous} forces equality
at every stage of the parabolic zero argument.  We record the resulting
rigidity before turning to the full classification.

\begin{theorem}[Rigidity of extremizers]
\label{thm:extremizer-structure}
Let \(p\ge5\) be prime, and suppose that
\[
 \mathcal S_p(\psi)=p^2-p+2.
\]
Choose a standard MUB in which \(\psi\) has minimum support and, using
Lemma~\ref{lem:standard-mub-transitivity}, normalize it to the
computational basis.  Put
\[
 k:=|\supp_{\mathcal B_\infty}(\psi)|.
\]
Then the following assertions hold.

\begin{enumerate}
\item The integer \(k\) is even.  Write \(k=2r\).

\item If \(k=2\), then every finite quadratic basis contains exactly one
zero coefficient.  Moreover, if the computational support is
\(\{u,v\}\), then, up to a nonzero scalar,
\[
 \psi=e_u-\omega^t e_v
\]
for some \(t\in\F_p\).

\item Suppose \(k\ge4\).  Then, after multiplying \(\psi\) by a nonzero
complex scalar, its nonzero computational coordinates belong to
\(K=\mathbb Q(\omega)\).  In the notation of
Lemma~\ref{lem:schrodinger-jet}, the first nonzero jet has order
\[
 n=r,
\]
and its leading pencil factors as
\begin{equation}
 Q_0(a,b)=B(a)\bigl(b+L(a)\bigr),
 \label{eq:extremizer-leading-factorization}
\end{equation}
where
\[
 \deg B=r-1,
 \qquad
 \deg L\le1,
\]
and \(B\) splits completely over \(\F_p\).

If the distinct roots of \(B\) are
\(\alpha_1,\ldots,\alpha_t\) with multiplicities
\(m_1,\ldots,m_t\), then
\begin{equation}
 m_1+\cdots+m_t=r-1.
 \label{eq:extremizer-partition}
\end{equation}
The \(F\)-fiber \(a=\alpha_i\) has exactly
\begin{equation}
 2m_i+1
 \label{eq:extremizer-exceptional-zeros}
\end{equation}
zeros.  Equivalently, putting \(A_i=-\alpha_i\), the physical finite basis
\(\mathcal B_{A_i}\) has exactly \(2m_i+1\) zero coefficients, while every
other physical finite basis has exactly one zero coefficient.

\item Since the computational basis has minimum support, each
\(m_i\) satisfies
\begin{equation}
 2m_i+1\le p-k.
 \label{eq:extremizer-minimum-support-constraint}
\end{equation}
Consequently, if \(k\ge4\), then necessarily \(k\le p-3\).
\end{enumerate}

Thus, up to permutation of the \(p+1\) basis labels, the zero-count
profile of an extremizer with minimum support \(k=2r\ge4\) has the form
\begin{equation}
 \boxed{
 \bigl(
 p-k,\,
 2m_1+1,\ldots,2m_t+1,\,
 \underbrace{1,\ldots,1}_{p-t\text{ entries}}
 \bigr),
 }
 \label{eq:extremizer-zero-profile}
\end{equation}
where \((m_1,\ldots,m_t)\) is a partition of \(r-1\) satisfying
\eqref{eq:extremizer-minimum-support-constraint}.
\end{theorem}

\begin{proof}
Choose the minimum-support basis and normalize it to
\(\mathcal B_\infty\), exactly as in the proof of
Theorem~\ref{thm:exact-mub-rigorous}.  Let
\[
 R=\supp_{\mathcal B_\infty}(\psi),
 \qquad |R|=k,
\]
and let \(N\) denote the total number of zero coefficients among the
\(p\) finite quadratic bases.  Then
\begin{equation}
 \mathcal S_p(\psi)=k+p^2-N.
 \label{eq:extremizer-S-vs-N}
\end{equation}
Since \(\mathcal S_p(\psi)=p^2-p+2\),
\begin{equation}
 N=p+k-2.
 \label{eq:extremizer-N}
\end{equation}

Theorem~\ref{thm:parabolic-zero-rigorous} gives
\[
 N\le p+2\left\lfloor\frac{k}{2}\right\rfloor-2.
\]
Comparing with \eqref{eq:extremizer-N} yields
\[
 k\le2\left\lfloor\frac{k}{2}\right\rfloor.
\]
The reverse inequality is automatic, hence equality holds and \(k\) is
even.  Write \(k=2r\).  This proves (1).

Assume first \(k=2\), and write
\[
 R=\{u,v\},
 \qquad
 \psi=c_u e_u+c_v e_v,
 \qquad c_uc_v\ne0.
\]
Equation \eqref{eq:extremizer-N} gives \(N=p\).  On the other hand,
for every fixed \(a\in\F_p\), the function
\[
 b\longmapsto
 c_u\omega^{a u^2+bu}+c_v\omega^{a v^2+bv}
\]
has at most one zero, because after division by the nonzero first term
a zero is equivalent to one equation in the nontrivial character
\(b\mapsto\omega^{b(v-u)}\).  Since there are \(p\) finite bases and
there are altogether \(p\) zeros, each finite basis has exactly one.
In particular, for \(a=0\) there is \(b_0\in\F_p\) such that
\[
 c_u\omega^{b_0u}+c_v\omega^{b_0v}=0.
\]
Therefore
\[
 \frac{c_v}{c_u}=-\omega^{b_0(u-v)}.
\]
After multiplication by \(c_u^{-1}\), this is the asserted form
\(e_u-\omega^t e_v\).  This proves (2).

Now suppose \(k=2r\ge4\).  We revisit the proof of
Theorem~\ref{thm:parabolic-zero-rigorous} and show that every inequality
used there must be an equality.

Let \(Z\) be the zero set of the corresponding parabolic sum, so
\(|Z|=N=p+2r-2\).  If the zero-row rank were at most \(k-2\), the proof
of Lemma~\ref{lem:complex-to-cyclotomic} would produce a nonzero
parabolic sum \(F'\) on an actual support \(R'\) with
\(|R'|\le k-1\), while preserving every zero in \(Z\).  If
\(|R'|=1\), then \(F'\) has no zeros, impossible.  If \(|R'|\ge2\),
Theorem~\ref{thm:parabolic-zero-rigorous} would give
\[
 |Z|
 \le
 p+2\left\lfloor\frac{|R'|}{2}\right\rfloor-2
 \le
 p+2(r-1)-2
 =p+2r-4,
\]
contrary to \(|Z|=p+2r-2\).  Hence the zero-row rank is exactly
\(k-1\).  Lemma~\ref{lem:complex-to-cyclotomic} therefore implies that,
after a nonzero scalar multiplication, the coefficient vector belongs
to \(K^R\).

We may consequently use the cyclotomic proof of
Lemma~\ref{lem:cyclotomic-zero-rigorous}.  Let \(n\) be the first
nonzero jet order.  That proof gives
\begin{equation}
 N\le p+2n-2
 \le p+2r-2.
 \label{eq:extremizer-chain-n}
\end{equation}
Since the two endpoints are equal to \(N\), both inequalities in
\eqref{eq:extremizer-chain-n} are equalities.  Thus
\begin{equation}
 n=r.
 \label{eq:extremizer-n-r}
\end{equation}

Write
\[
 Q_0(a,b)=A(a)+bB(a),
 \qquad
 \deg A\le r,
 \quad
 \deg B\le r-1.
\]
We first show that \(B\not\equiv0\).  If \(B\equiv0\), Case 2 in the
proof of Lemma~\ref{lem:cyclotomic-zero-rigorous} gives, away from the
single separately treated boundary configuration, the stronger bound
\[
 N\le3n=3r.
\]
But
\[
 3r<p+2r-2,
\]
because \(r\le(p-1)/2\) and \(p\ge5\).  In the boundary configuration
that proof gives the still stronger estimate \(N\le p-2\).  Both
contradict \(N=p+2r-2\).  Hence \(B\not\equiv0\).

Let \(D=\gcd(A,B)\), chosen monic.  Case 1 of the same proof gives the
chain
\begin{equation}
 \begin{aligned}
 N
 &\le
 p+2\sum_{\alpha\in Z_{\F_p}(D)}m_\alpha\\
 &\le p+2\deg D\\
 &\le p+2\deg B\\
 &\le p+2r-2,
 \end{aligned}
 \label{eq:extremizer-equality-chain}
\end{equation}
where
\[
 m_\alpha=
 \min\{\operatorname{ord}_\alpha A,
        \operatorname{ord}_\alpha B\}.
\]
Again the first and last quantities in
\eqref{eq:extremizer-equality-chain} are both equal to \(N\), so every
inequality in that chain is an equality.  In particular,
\begin{equation}
 \deg B=r-1,
 \qquad
 \deg D=\deg B.
 \label{eq:extremizer-degrees}
\end{equation}
Since \(D\mid B\) and the two polynomials have the same degree,
\(D\) differs from \(B\) only by a nonzero scalar.  Thus \(B\mid A\).
Because \(\deg A\le r\) and \(\deg B=r-1\), there is a polynomial
\(L\) of degree at most one such that
\[
 A=BL.
\]
This proves the factorization
\eqref{eq:extremizer-leading-factorization}.

Equality in
\[
 \sum_{\alpha\in Z_{\F_p}(D)}m_\alpha\le\deg D
\]
forces every root of \(D\), with its full multiplicity, to lie in
\(\F_p\).  Since \(D\) is a scalar multiple of \(B\), the polynomial
\(B\) splits completely over \(\F_p\).  If its distinct roots are
\(\alpha_1,\ldots,\alpha_t\) with multiplicities
\(m_1,\ldots,m_t\), then
\[
 m_1+\cdots+m_t=\deg B=r-1,
\]
which is \eqref{eq:extremizer-partition}.

It remains to prove the exact fiber counts.  In Case 1 of the zero-bound
proof, a fiber with \(B(\alpha)\ne0\) contributes at most one exact
zero, while a root \(\alpha_i\) of multiplicity \(m_i\) contributes at
most \(2m_i+1\) exact zeros.  Summing these individual upper bounds now
gives exactly
\[
 p-t+\sum_{i=1}^t(2m_i+1)
 =p+2\sum_{i=1}^t m_i
 =p+2r-2
 =N.
\]
Since the sum of the fiberwise upper bounds equals the actual total,
every fiberwise upper bound must itself be attained.  Hence each
exceptional fiber \(a=\alpha_i\) contains exactly \(2m_i+1\) zeros and
every nonexceptional finite fiber contains exactly one.  This proves
(3).

Finally, the computational basis was chosen to have minimum support
\(k\).  Therefore every finite basis has support at least \(k\), or,
equivalently, has at most \(p-k\) zero coefficients.  Applying this to the physical exceptional basis
\(\mathcal B_{A_i}\), where \(A_i=-\alpha_i\), gives
\[
 2m_i+1\le p-k,
\]
which is \eqref{eq:extremizer-minimum-support-constraint}.  If
\(k\ge4\), then \(r-1\ge1\), so at least one exceptional root exists
and its multiplicity is at least one.  Hence \(p-k\ge3\), i.e.
\(k\le p-3\).  This proves (4) and completes the proof.
\end{proof}

\section{Classification and enumeration of extremizers}
\label{sec:complete-extremizers}

We now turn the equality constraints into a complete parametrization of the
extremizing rays in every prime dimension \(p\ge5\).

Put
\[
 \mathcal P:=\F_p^\times/\{\pm1\},
 \qquad
 \kappa:=|\mathcal P|=\frac{p-1}{2}.
\]
We write \([t]\) for the class of \(t\in\F_p^\times\).

\begin{lemma}[Universal constraint matrix]
\label{lem:complete-constraint-matrix}
Let \(r\ge2\).  Choose
\[
 \mathcal R=\{[\rho_1],\ldots,[\rho_r]\}
 \in\binom{\mathcal P}{r}
\]
and
\[
 \mathcal C\in
 \binom{\F_p\times\mathcal P}{r-1}.
\]
Choose representatives \(\rho_j\in\F_p^\times\), and for each
\((\alpha,[e])\in\mathcal C\) choose \(e\in\F_p^\times\).  Define
\begin{equation}
 M_{\mathcal C,\mathcal R}
 =
 \left(
 \omega^{\alpha\rho_j^2}
 \bigl(\omega^{e\rho_j}-\omega^{-e\rho_j}\bigr)
 \right)_{
 (\alpha,[e])\in\mathcal C,\ 1\le j\le r}.
 \label{eq:complete-constraint-matrix}
\end{equation}
Then:

\begin{enumerate}
\item
\(\operatorname{rank}M_{\mathcal C,\mathcal R}=r-1\);
\item every maximal minor is nonzero;
\item \(\ker M_{\mathcal C,\mathcal R}\) is generated by a vector
\[
 q=(q_1,\ldots,q_r)
\]
with \(q_j\ne0\) for every \(j\).
\end{enumerate}

Changing a representative \(e\) to \(-e\) multiplies the corresponding
row by \(-1\) and leaves the kernel unchanged.  Changing one radius
representative \(\rho_j\) to \(-\rho_j\) multiplies the corresponding column
by \(-1\).  More generally, if \(S\) is the resulting diagonal sign matrix,
then
\[
 M'_{\mathcal C,\mathcal R}=M_{\mathcal C,\mathcal R}S,
 \qquad
 \ker M'_{\mathcal C,\mathcal R}=S^{-1}\ker M_{\mathcal C,\mathcal R}.
\]
Thus the kernel ray is covariant, rather than literally invariant, under
changes of the radius representatives.  In the state construction of
Theorem~\ref{thm:complete-extremizer-classification}, this column sign is
compensated by exchanging the two support labels \(c\pm\rho_j\); hence the
resulting projective state is independent of all representative choices.
\end{lemma}

\begin{proof}
Suppose a maximal minor vanishes.  Then there is a nonzero vector
\[
 u=(u_1,\ldots,u_r)
\]
supported on at most \(r-1\) columns such that
\[
 M_{\mathcal C,\mathcal R}u=0.
\]
Let \(s\le r-1\) be the number of nonzero coordinates of \(u\), and
form the \(b\)-odd parabolic sum
\[
 G(a,b)
 =
 \sum_{j:u_j\ne0}
 u_j\omega^{a\rho_j^2}
 \bigl(\omega^{b\rho_j}-\omega^{-b\rho_j}\bigr).
\]
Its actual frequency support on the parabola has size \(2s\).

For every \(a\in\F_p\),
\[
 G(a,0)=0.
\]
Thus \(G\) has \(p\) central zeros.  For every
\((\alpha,[e])\in\mathcal C\), the row relation gives
\[
 G(\alpha,e)=0,
\]
and oddness in \(b\) also gives
\[
 G(\alpha,-e)=0.
\]
Because the \(r-1\) elements of \(\mathcal C\) are distinct as pairs
\((\alpha,[e])\), these are \(2(r-1)\) distinct noncentral zeros.
Hence
\[
 \#Z(G)\ge p+2r-2.
\]

Theorem~\ref{thm:parabolic-zero-rigorous}, applied to the actual support
size \(2s\), gives
\[
 \#Z(G)\le p+2s-2\le p+2r-4,
\]
a contradiction.  Hence every maximal minor is nonzero.  The remaining
assertions follow from the cofactor formula for the one-dimensional
kernel.
\end{proof}

\begin{theorem}[Classification of all extremizers]
\label{thm:complete-extremizer-classification}
Let \(p\ge5\) be prime and let
\[
 \mathcal S_p(\psi)=p^2-p+2.
\]
Choose one standard MUB in which \(\psi\) has minimum support, fix a
standard-MUB-permuting unitary that sends it to the computational basis,
and work with the resulting normalized projective state.

Then the minimum support is
\begin{equation}
 k=2r
 \qquad\text{for some}\qquad
 1\le r\le\frac{p-3}{2}.
 \label{eq:allowed-minimum-supports}
\end{equation}
There are unique data
\begin{equation}
 c,\nu\in\F_p,
 \qquad
 \mathcal R\in\binom{\mathcal P}{r},
 \qquad
 \mathcal C\in
 \binom{\F_p\times\mathcal P}{r-1},
 \label{eq:complete-extremizer-data}
\end{equation}
such that, with
\[
 m_\alpha
 :=
 \#\{[e]\in\mathcal P:(\alpha,[e])\in\mathcal C\},
\]
one has the capacity inequalities
\begin{equation}
 m_\alpha
 \le
 M_{p,r}:=
 \frac{p-2r-1}{2}
 \qquad(\alpha\in\F_p).
 \label{eq:complete-capacity}
\end{equation}

For \(r=1\), the set \(\mathcal C\) is empty and we use the convention
\(q_1=1\).  For \(r\ge2\), write
\[
 \mathcal R=\{[\rho_1],\ldots,[\rho_r]\}
\]
and choose representatives \(\rho_j\) and representatives \(e\) for the
classes occurring in \(\mathcal C\), as in
Lemma~\ref{lem:complete-constraint-matrix}.  Let
\[
 q=(q_1,\ldots,q_r)
\]
generate the corresponding kernel of
\(M_{\mathcal C,\mathcal R}\).  After multiplication of
\(\psi\) by one nonzero scalar,
\begin{equation}
 \boxed{
 \begin{aligned}
 \psi_{c+\rho_j}
 &=
 q_j\,\omega^{-\nu(c+\rho_j)},\\
 \psi_{c-\rho_j}
 &=
 -q_j\,\omega^{-\nu(c-\rho_j)}
 \qquad(1\le j\le r),
 \end{aligned}}
 \label{eq:complete-state-formula}
\end{equation}
and all remaining computational coordinates vanish.  This projective state
is independent of the representative choices: changing an \(e\)-representative
only changes a row sign, while replacing \(\rho_j\) by \(-\rho_j\) changes
\(q_j\) to \(-q_j\) and simultaneously exchanges the two support labels
\(c+\rho_j\) and \(c-\rho_j\), leaving the coordinates in
\eqref{eq:complete-state-formula} unchanged after that identification.

For the parabolic sum
\[
 F(a,b)=\sum_x\psi_x\omega^{a x^2+b x},
\]
the complete zero set in \(F\)-coordinates is
\begin{equation}
 \boxed{
 \begin{aligned}
 b&=\nu-2ac
 &&(a\in\F_p),\\
 b&=\nu-2\alpha c\pm e
 &&((\alpha,[e])\in\mathcal C).
 \end{aligned}}
 \label{eq:complete-zero-set}
\end{equation}
Equivalently, in the physical MUB labels \((A,\beta)\), the zero
coefficients occur exactly at
\begin{equation}
 \boxed{
 \begin{aligned}
 \beta&=-\nu-2Ac
 &&(A\in\F_p),\\
 A&=-\alpha,\qquad
 \beta=-\nu+2\alpha c\pm e
 &&((\alpha,[e])\in\mathcal C).
 \end{aligned}}
 \label{eq:complete-physical-zero-set}
\end{equation}
Thus the \(F\)-fiber \(a=\alpha\) has exactly
\[
 1+2m_\alpha
\]
zeros; equivalently, the physical basis \(\mathcal B_A\) contains exactly
\(1+2m_{-A}\) zero coefficients.

Conversely, every choice of data
\eqref{eq:complete-extremizer-data} satisfying
\eqref{eq:complete-capacity} gives, through
\eqref{eq:complete-state-formula}, a sharp state whose minimum support
is \(2r\).  Once the MUB-permuting normalization above has been fixed,
the data are unique for the resulting normalized projective state.
\end{theorem}

\begin{proof}
We first treat minimum support two.  Suppose
\[
 R=\{u,v\},
 \qquad u\ne v.
\]
The computational support has size two.  By the pairwise support
inequality, every one of the other \(p\) MUBs has support at least \(p-1\).
Since the total support equals
\[
 p^2-p+2=2+p(p-1),
\]
the sum of those \(p\) remaining support sizes is exactly \(p(p-1)\).
Consequently every one of them has support exactly \(p-1\).  In particular
the Fourier basis has a zero.  Therefore, for some \(b_0\in\F_p\),
\[
 \psi_u\omega^{b_0u}
 +
 \psi_v\omega^{b_0v}
 =0.
\]
Thus
\[
 \frac{\psi_v}{\psi_u}
 =
 -\omega^{b_0(u-v)}.
\]
Putting
\[
 c=\frac{u+v}{2},
 \qquad
 \rho=\frac{u-v}{2},
\]
and absorbing one nonzero scalar, this is exactly
\eqref{eq:complete-state-formula} with \(r=1\), empty
\(\mathcal C\), and a suitable \(\nu\).  Direct substitution gives the
zero line
\[
 b=\nu-2ac.
\]
This is the asserted classification for \(r=1\).

Assume from now on that the minimum support is at least four.  By the
parity statement in Theorem~\ref{thm:extremizer-structure}, it equals
\(2r\) for some \(r\ge2\).  The same theorem gives
\begin{equation}
 Q_0(a,b)=B(a)\bigl(b+L(a)\bigr),
 \qquad
 \deg B=r-1,
 \qquad
 \deg L\le1,
 \label{eq:complete-leading-pencil}
\end{equation}
where \(B\) splits completely over \(\F_p\).  If its distinct roots are
\[
 \alpha_1,\ldots,\alpha_t
\]
with multiplicities
\[
 m_1,\ldots,m_t,
\]
then
\begin{equation}
 m_1+\cdots+m_t=r-1,
 \qquad
 t\le r-1,
 \label{eq:complete-partition}
\end{equation}
the fiber \(a=\alpha_i\) contains exactly \(2m_i+1\) exact zeros, and
every nonexceptional finite fiber contains exactly one exact zero.

Write
\[
 -L(a)=\mu a+\nu.
\]
Every nonexceptional fiber therefore has its unique zero at
\[
 b=\mu a+\nu.
\]
Put
\[
 d_x:=\psi_x\omega^{\nu x},
 \qquad
 h(x):=x^2+\mu x,
\]
and
\[
 g(t):=
 \sum_{\substack{x\in R\\h(x)=t}}d_x.
\]
For every nonexceptional \(a\),
\[
 \widehat g(a)=0.
\]
Hence
\begin{equation}
 |\supp\widehat g|
 \le t
 \le r-1.
 \label{eq:complete-ghat-support}
\end{equation}

Let
\[
 c=-\frac{\mu}{2}.
\]
Then
\[
 h(x)=(x-c)^2-c^2.
\]
The set of squares in \(\F_p\), including \(0\), has
\((p+1)/2\) elements.  Consequently
\begin{equation}
 |\supp g|
 \le|\operatorname{im}h|
 =\frac{p+1}{2}.
 \label{eq:complete-g-support}
\end{equation}
If \(g\ne0\), Tao's prime Fourier uncertainty principle gives
\[
 p+1
 \le
 |\supp g|+|\supp\widehat g|.
\]
But the minimum support is at most \(p-1\), so
\[
 r\le\frac{p-1}{2},
\]
and therefore
\[
 |\supp g|+|\supp\widehat g|
 \le
 \frac{p+1}{2}+(r-1)
 \le
 \frac{p+1}{2}+\frac{p-3}{2}
 =p-1,
\]
a contradiction.  Hence
\begin{equation}
 g\equiv0.
 \label{eq:complete-g-zero}
\end{equation}

The fibers of \(h\) are the orbits of the involution
\[
 x\longmapsto2c-x.
\]
The fixed fiber is \(\{c\}\).  If \(c\in R\), then
\(g(h(c))=d_c\ne0\), contradicting \eqref{eq:complete-g-zero}; hence
\(c\notin R\).  Every other fiber has the form
\(\{c-\rho,c+\rho\}\).  If such a fiber met \(R\) in exactly one point,
then its contribution to \(g\) would be a single nonzero coefficient,
again contradicting \(g\equiv0\).  Thus every occupied nonfixed fiber is
occupied at both points, and on that fiber the equation \(g=0\) reads
\[
 d_{c-\rho}+d_{c+\rho}=0.
\]
Since \(|R|=2r\), there are exactly \(r\) occupied two-point fibers.
Therefore
\[
 R=
 \{c\pm\rho_1,\ldots,c\pm\rho_r\}
\]
for distinct radius classes
\[
 \mathcal R=
 \{[\rho_1],\ldots,[\rho_r]\}\subset\mathcal P,
\]
and
\begin{equation}
 d_{c-\rho_j}=-d_{c+\rho_j}
 \qquad(1\le j\le r).
 \label{eq:complete-antisymmetry}
\end{equation}
Put
\[
 q_j:=d_{c+\rho_j}.
\]
This gives \eqref{eq:complete-state-formula}.

The antisymmetry produces one central zero in every \(F\)-fiber:
\[
 F(a,\nu-2ac)=0
 \qquad(a\in\F_p).
\]
Equivalently, every physical finite basis \(\mathcal B_A\) has the
central zero at \(\beta=-\nu-2Ac\).
For
\[
 H_a(t):=F(a,\nu-2ac+t),
\]
the contribution of the pair \(c\pm\rho_j\), using
\eqref{eq:complete-state-formula}, is
\[
 q_j\omega^{-ac^2+a\rho_j^2+ct}
 \bigl(\omega^{t\rho_j}-\omega^{-t\rho_j}\bigr).
\]
Replacing \(t\) by \(-t\) changes the parenthesis by a factor \(-1\)
and changes \(\omega^{ct}\) to \(\omega^{-ct}\).  Summing over \(j\)
gives
\begin{equation}
 H_a(-t)
 =
 -\omega^{-2ct}H_a(t).
 \label{eq:complete-fiber-symmetry}
\end{equation}
Thus every fiber zero set is symmetric about \(t=0\).  The
exceptional fiber \(a=\alpha_i\), which has \(2m_i+1\) zeros, therefore
has the form
\[
 t=0,\quad
 t=\pm e_{i,1},\ldots,\pm e_{i,m_i},
\]
where the classes \([e_{i,\ell}]\) are distinct.  Define
\[
 \mathcal C
 :=
 \{(\alpha_i,[e_{i,\ell}]):
 1\le i\le t,\ 1\le\ell\le m_i\}.
\]
Then
\[
 |\mathcal C|
 =
 \sum_i m_i
 =
 r-1.
\]
The noncentral zero equations are exactly
\[
 M_{\mathcal C,\mathcal R}q=0.
\]
Lemma~\ref{lem:complete-constraint-matrix} shows that this kernel ray is
unique and that all \(q_j\ne0\).

Since the computational basis was chosen to have minimum support
\(2r\), the physical basis \(\mathcal B_A\), with \(A=-\alpha\), corresponding
to every exceptional \(F\)-fiber \(a=\alpha\) has support at least \(2r\):
\[
 p-(2m_\alpha+1)\ge2r.
\]
This is exactly \eqref{eq:complete-capacity}.  Since
\(|\mathcal C|=r-1>0\), some \(m_\alpha\ge1\).  Hence
\[
 M_{p,r}\ge1,
\]
which implies
\[
 2r\le p-3.
\]
Together with the already treated case \(r=1\), this proves
\eqref{eq:allowed-minimum-supports}.

The data are uniquely recoverable from the resulting normalized
projective state.  Indeed,
\[
 c=\frac{1}{2r}\sum_{x\in R}x
\]
because \(2r<p\); the support determines the radius classes
\(\mathcal R\).  Each \(F\)-fiber zero set is a nonempty proper subset of \(\F_p\)
invariant under reflection about its displayed \(b\)-center, and that center
is unique: invariance under reflections about two distinct centers would
imply invariance under a nonzero translation, hence under all translations
of the prime cyclic group, forcing the zero set to be all of \(\F_p\), which
is impossible for a nonzero state.  Thus the centers of the symmetric
\(F\)-fiber zero sets determine the affine line \(b=\nu-2ac\), hence
\(\nu\); and the noncentral symmetric zero pairs determine \(\mathcal C\).
By \eqref{eq:F-to-physical-labels}, the equivalent physical center in
\(\mathcal B_A\) is \(\beta=-\nu-2Ac\).

Conversely, choose data satisfying
\eqref{eq:complete-extremizer-data} and
\eqref{eq:complete-capacity}.  For \(r=1\), the formula gives a
two-point state and direct substitution shows one zero in every finite
basis, so the state is sharp.

Let \(r\ge2\).  By
Lemma~\ref{lem:complete-constraint-matrix}, the constraint matrix has a
unique kernel ray generated by \(q\) with no zero coordinate.  Define
\(\psi\) by \eqref{eq:complete-state-formula}.  Its computational
support has size \(2r\).

Centered antisymmetry gives the \(p\) central zeros
\[
 b=\nu-2ac.
\]
Each constraint \((\alpha,[e])\in\mathcal C\) gives the two additional
zeros
\[
 b=\nu-2\alpha c\pm e.
\]
Thus
\[
 \#Z(F)\ge
 p+2(r-1)
 =
 p+2r-2.
\]
The parabolic zero theorem gives the reverse inequality for support
\(2r\):
\[
 \#Z(F)\le p+2r-2.
\]
Hence the displayed points are exactly all zeros of \(F\), and under
\eqref{eq:F-to-physical-labels} they correspond bijectively to all zero
coefficients in the finite MUBs.  Therefore
\[
 \mathcal S_p(\psi)
 =
 2r+p^2-(p+2r-2)
 =
 p^2-p+2.
\]
The \(F\)-fiber \(a=\alpha\) has exactly \(1+2m_\alpha\) zeros, so the
corresponding physical basis \(\mathcal B_A\), where \(A=-\alpha\), has support
\[
 p-(1+2m_\alpha)\ge2r
\]
by \eqref{eq:complete-capacity}.  Every physical finite basis corresponding
to a nonexceptional \(F\)-fiber has support \(p-1>2r\).  Therefore the chosen computational basis is indeed
a minimum-support basis.  This proves the converse.
\end{proof}

\begin{corollary}[Possible minimum supports]
\label{cor:possible-minimum-supports}
For \(p\ge5\), the minimum support of a sharp state is exactly one of
\[
 \boxed{
 2,4,6,\ldots,p-3.
 }
\]
Every value in this list occurs.
\end{corollary}

\begin{proof}
Necessity follows from
Theorem~\ref{thm:complete-extremizer-classification}.  Conversely, let
\(2r\le p-3\).  Then
\[
 M_{p,r}=\frac{p-2r-1}{2}\ge1.
\]
Choose \(r\) radius classes and choose the \(r-1\) constraints in
distinct \(a\)-fibers, so every occupancy \(m_\alpha\) is at most one.
The capacity condition holds, and
Theorem~\ref{thm:complete-extremizer-classification} produces a sharp
state of minimum support \(2r\).
\end{proof}

\begin{corollary}[Sharpness at every even support]
\label{cor:even-parabolic-sharpness}
For every prime $p\ge5$ and every integer \(r\) with
\[
 1\le r\le\frac{p-1}{2},
\]
there exist $R\subset\mathbb F_p$ with $|R|=2r$ and nonzero complex
coefficients $(c_x)_{x\in R}$ such that
\[
 \#Z\!\left(\sum_{x\in R}c_x\omega^{ax^2+bx}\right)=p+2r-2.
\]
Hence the even-support bound in Theorem~\ref{thm:intro-parabolic} is sharp
throughout its full range.
\end{corollary}

\begin{proof}
First suppose \(2r\le p-3\).  By
Corollary~\ref{cor:possible-minimum-supports}, choose a sharp MUB state whose
minimum support is \(2r\), and normalize one minimum-support basis to
\(\mathcal B_\infty\).  If \(N\) is the number of zero coefficients among
the \(p\) finite quadratic bases, then
\[
 p^2-p+2=2r+p^2-N,
\]
so \(N=p+2r-2\).  By \eqref{eq:mub-parabola-label-map} and the bijection
\((A,\beta)\mapsto(-A,-\beta)\), these are exactly the zeros of the
associated parabolic sum.

It remains to treat the endpoint \(2r=p-1\).  Take
\[
 R=\F_p^\times,
 \qquad
 c_x=1-\omega^{-x}\quad(x\in R).
\]
All these coefficients are nonzero.  Since the coefficient at \(x=0\)
would be zero, we may extend the sum to all \(x\in\F_p\) and write
\[
 F(a,b)
 =
 \sum_{x\in\F_p}\omega^{ax^2+bx}
 -
 \sum_{x\in\F_p}\omega^{ax^2+(b-1)x}.
\]
For \(a=0\), character orthogonality gives
\[
 F(0,b)=p\,\mathbf 1_{\{b=0\}}-p\,\mathbf 1_{\{b=1\}},
\]
so the fiber \(a=0\) contains exactly \(p-2\) zeros.  For \(a\ne0\),
completing the square gives
\[
 \sum_{x\in\F_p}\omega^{ax^2+bx}
 =
 G(a)\omega^{-b^2/(4a)},
 \qquad
 G(a):=\sum_{x\in\F_p}\omega^{ax^2}\ne0.
\]
Hence \(F(a,b)=0\) if and only if
\[
 b^2=(b-1)^2,
\]
that is, \(b=\tfrac12\).  Thus every one of the \(p-1\) nonzero
\(a\)-fibers contributes exactly one zero, and
\[
 |Z(F)|=(p-2)+(p-1)=2p-3=p+(p-1)-2.
\]
This is equality in the even-support bound for \(2r=p-1\).
\end{proof}

\begin{corollary}[Enumeration of extremizers]
\label{cor:complete-extremizer-count}
For
\[
 1\le r\le\frac{p-3}{2},
\]
put
\[
 M=M_{p,r}=\frac{p-2r-1}{2},
 \qquad
 \kappa=\frac{p-1}{2},
\]
and define
\begin{equation}
 P_{p,r}(z,u)
 :=
 \sum_{j=0}^{M-1}\binom{\kappa}{j} z^j
 +
 u\binom{\kappa}{M}z^M.
 \label{eq:counting-polynomial}
\end{equation}
Then the number \(N_{2r}(p)\) of sharp projective rays whose minimum
support is \(2r\) is
\begin{equation}
 \boxed{
 N_{2r}(p)
 =
 (p+1)p^2\binom{\kappa}{r}
 \,[z^{r-1}]
 \int_0^1
 P_{p,r}(z,u)^p\,du.
 }
 \label{eq:complete-count-formula}
\end{equation}
The integral is coefficientwise in the polynomial variable \(u\).
Consequently the total number of sharp rays is
\begin{equation}
 N_{\mathrm{sharp}}(p)
 =
 \sum_{r=1}^{(p-3)/2}N_{2r}(p).
 \label{eq:total-sharp-count}
\end{equation}
\end{corollary}

\begin{proof}
Fix one standard MUB and designate it as the chosen minimum-support
basis.  By the classification theorem, a normalized extremizer of
minimum support \(2r\) is determined uniquely by:

\begin{enumerate}
\item a center \(c\in\F_p\);
\item a parameter \(\nu\in\F_p\);
\item \(r\) radius classes chosen from the \(\kappa=(p-1)/2\) classes in
\(\mathcal P\);
\item a constraint set
\[
 \mathcal C\subset\F_p\times\mathcal P,
 \qquad
 |\mathcal C|=r-1,
\]
whose occupancy in each \(a\)-fiber is at most \(M\).
\end{enumerate}

For a fixed \(a\), choosing exactly \(j\) constraints in that fiber can
be done in
\[
 \binom{\kappa}{j}
\]
ways.  The polynomial
\[
 P_{p,r}(z,u)
\]
records this choice, with \(z\) marking the number of constraints and
\(u\) marking a saturated fiber, that is, a fiber with occupancy
exactly \(M\).  Hence
\[
 [z^{r-1}u^s]P_{p,r}(z,u)^p
\]
counts the constraint sets with exactly \(s\) saturated finite fibers.

For a finite fiber with occupancy \(m_\alpha\), the classification
theorem gives support
\[
 p-(1+2m_\alpha).
\]
By definition
\[
 M=\frac{p-2r-1}{2},
\]
so this support equals \(2r\) if and only if \(m_\alpha=M\); if
\(m_\alpha<M\), it is strictly larger than \(2r\).  Thus the saturated \(F\)-fibers are exactly those whose corresponding
physical bases \(\mathcal B_A\), with \(A=-\alpha\), also have minimum support.  A
constraint set with exactly \(s\) saturated fibers therefore produces exactly
\[
 1+s
\]
minimum-support bases: the chosen computational basis and those \(s\) finite
bases.  Consequently a count in which one designates a minimum-support basis
counts the same projective ray exactly \(1+s\) times.

Fix once and for all, for each standard MUB, one MUB-permuting unitary
that sends it to the computational basis.  This gives a bijection between
rays with that designated minimum-support basis and the normalized rays
counted above.  Hence there are \(p+1\) choices for the designated basis,
and the remaining choices contribute
\[
 p^2\binom{\kappa}{r}.
\]
Therefore
\[
 N_{2r}(p)
 =
 (p+1)p^2\binom{\kappa}{r}
 \sum_{s\ge0}
 \frac{1}{s+1}
 [z^{r-1}u^s]P_{p,r}(z,u)^p.
\]
Since
\[
 \frac1{s+1}=\int_0^1u^s\,du,
\]
this is exactly \eqref{eq:complete-count-formula}.
\end{proof}

\begin{corollary}[Small-prime counts]
\label{cor:small-prime-sharp-counts}
The enumeration formula gives
\[
 N_{\mathrm{sharp}}(5)=300,
\]
as also follows directly from the two-point case, and
\[
 N_{\mathrm{sharp}}(7)=13524.
\]
In dimension eleven, the entries below are indexed by the minimum support
size \(k\):
\begin{center}
\begin{tabular}{c@{\qquad}rrrr}
\toprule
minimum support \(k\) & 2 & 4 & 6 & 8 \\
\midrule
\(N_k(11)\) & 7260 & 798600 & 20763600 & 37434375 \\
\bottomrule
\end{tabular}
\end{center}
and hence
\[
 \boxed{
 N_{\mathrm{sharp}}(11)=59003835.
 }
\]
These displayed values were independently checked by exact integer/rational
coefficient expansion of the generating polynomial in
\eqref{eq:counting-polynomial} and coefficient extraction in
\eqref{eq:complete-count-formula}; no floating-point zero test is involved.
A short verification script is included with the source files.
\end{corollary}

\begin{remark}[Geometry of equality]
\label{rem:complete-equality-geometry}
The classification shows that every sharp state is governed by two
reflection structures.

First, its minimum-support coordinates occur in pairs
\[
 c\pm\rho_j
\]
and the dephased coefficients are odd under reflection about \(c\).
Second, in \(F\)-coordinates the fiberwise zero sets are symmetric about
the affine line
\[
 b=\nu-2ac.
\]
Equivalently, in the physical finite MUB labels the centers lie on
\[
 \beta=-\nu-2Ac.
\]
The noncentral zero pairs form an arbitrary admissible set
\[
 \mathcal C\subset\F_p\times\mathcal P
\]
subject only to the capacity condition
\eqref{eq:complete-capacity}; the universal nonvanishing of the
constraint-matrix maximal minors then produces the unique extremizing
ray.
\end{remark}


\section{Coding and matroid consequences}
\label{sec:code-matroid}

Let
\[
 \mathcal V
 =
 \bigcup_{a\in\F_p\cup\{\infty\}}\mathcal B_a
\]
be the \(p(p+1)\)-element vector configuration formed by the standard
complete MUB, with one chosen unit representative for every projective
point.  Linear-dependence and spark questions for structured finite frames
and MUB-related configurations have been studied from several directions
\cite{CasazzaKutyniok,BodmannPaulsen,AlexeevCahillMixon,Blanchfield,
JafarpourDuarteCalderbank}.  Closely related frame-robustness questions under
erasures, including kernel-sparsity and structured full-spark constructions,
appear in \cite{LiFrameErasures,LiTNS2026}.  The invariants below are naturally
dual: they measure minimum Hamming weight in
the analysis range, equivalently maximum hyperplane sections and minimum
cocircuits of this particular complete-MUB configuration.

Define the analysis code
\begin{equation}
 \mathcal C_p
 :=
 \left\{
 \bigl(\langle v,\psi\rangle\bigr)_{v\in\mathcal V}:
 \psi\in\C^p
 \right\}
 \subset\C^{p(p+1)}.
 \label{eq:mub-analysis-code}
\end{equation}
This is a complex linear code of length \(p(p+1)\); we use standard coding
terminology as in \cite{HuffmanPless}.  The analysis map is injective, because
its coordinates coming from the computational basis are precisely the
computational coordinates of \(\psi\) (up to the fixed inner-product
convention).  Hence \(\dim_\C\mathcal C_p=p\).

\begin{corollary}[Minimum distance of the complete-MUB analysis code]
\label{cor:mub-code-distance}
For every prime \(p\ge5\),
\begin{equation}
 \boxed{
 d_{\min}(\mathcal C_p)=p^2-p+2.
 }
 \label{eq:mub-code-distance}
\end{equation}
Moreover the projective minimum-weight codewords are classified
completely by
Theorem~\ref{thm:complete-extremizer-classification}.
\end{corollary}

\begin{proof}
For
\[
 c(\psi)
 :=
 \bigl(\langle v,\psi\rangle\bigr)_{v\in\mathcal V},
\]
its Hamming weight is exactly the total support of \(\psi\) in the
\(p+1\) bases:
\[
 \operatorname{wt}_H(c(\psi))
 =
 \sum_{a\in\F_p\cup\{\infty\}}
 s_{\mathcal B_a}(\psi)
 =
 \mathcal S_p(\psi).
\]
Theorem~\ref{thm:exact-mub-rigorous} therefore gives
\[
 d_{\min}(\mathcal C_p)=p^2-p+2.
\]
The equality cases are precisely the sharp states classified in
Theorem~\ref{thm:complete-extremizer-classification}.
\end{proof}

Let \(M_p\) be the complex vector matroid represented by the columns
\(\mathcal V\); for standard matroid terminology and duality we refer to
\cite{Oxley}.  Its ground set has size
\[
 |\mathcal V|=p(p+1)
\]
and rank \(p\).

\begin{corollary}[Maximum hyperplane sections and cogirth]
\label{cor:mub-matroid-cogirth}
For every prime \(p\ge5\), the largest number of vectors of
\(\mathcal V\) contained in a complex hyperplane of \(\C^p\) is
\begin{equation}
 \boxed{2p-2.}
 \label{eq:mub-max-hyperplane}
\end{equation}
Equivalently, the cogirth of the represented matroid \(M_p\) is
\begin{equation}
 \boxed{
 g^*(M_p)=p^2-p+2.
 }
 \label{eq:mub-cogirth}
\end{equation}
The maximum hyperplane sections, equivalently the minimum cocircuits,
are classified by
Theorem~\ref{thm:complete-extremizer-classification}.
\end{corollary}

\begin{proof}
Every complex hyperplane can be written
\[
 H_\psi
 =
 \{x\in\C^p:\langle x,\psi\rangle=0\}
\]
for some nonzero \(\psi\).  Hence
\[
 |\mathcal V\cap H_\psi|
 =
 p(p+1)-\mathcal S_p(\psi).
\]
The exact support theorem gives
\[
 |\mathcal V\cap H_\psi|
 \le
 p(p+1)-(p^2-p+2)
 =
 2p-2.
\]
Sharp states attain equality, so
\eqref{eq:mub-max-hyperplane} is exact.

For a represented matroid, a cocircuit is the complement in the ground
set of a matroid hyperplane.  Every matroid hyperplane is of the form
\(\mathcal V\cap H\) for the linear hyperplane \(H\) spanned by that
rank-\((p-1)\) flat, so its cardinality is at most \(2p-2\).
Conversely, let \(H_\psi\) be a linear hyperplane attaining the maximum
intersection size \(2p-2\), and put \(X=\mathcal V\cap H_\psi\).  We first
show that \(X\) has matroid rank \(p-1\).  If its rank were at most \(p-2\),
then, because \(\mathcal V\) spans \(\C^p\), one could choose
\(v\in\mathcal V\setminus\operatorname{span}(X)\).  The space
\(\operatorname{span}(X\cup\{v\})\) would have dimension at most \(p-1\)
and could be extended to a complex linear hyperplane \(H'\).  Then
\(H'\cap\mathcal V\) would contain the \(2p-2\) elements of \(X\) together
with \(v\), contradicting the maximal linear-hyperplane intersection bound.
Thus \(\operatorname{rank}_{M_p}X=p-1\).

Moreover
\[
 \operatorname{cl}_{M_p}(X)
 =\mathcal V\cap\operatorname{span}(X).
\]
Since \(\operatorname{span}(X)\subseteq H_\psi\) and both spaces now have
dimension \(p-1\), they are equal.  Hence
\(\operatorname{cl}_{M_p}(X)=\mathcal V\cap H_\psi=X\), so \(X\) itself is
a matroid hyperplane of cardinality \(2p-2\).  Therefore the maximum
matroid-hyperplane size is exactly \(2p-2\), and the minimum cocircuit size is
\[
 p(p+1)-(2p-2)
 =
 p^2-p+2.
\]
The maximum sections and minimum cocircuits have the same ground-set
intersections, so the classification again follows from the equality
theorem.
\end{proof}

\begin{remark}
The parabolic zero theorem may therefore be read simultaneously as a
sharp support uncertainty principle, a minimum-distance theorem for a
highly structured cyclotomic code, and a cogirth theorem for the
complete-MUB vector matroid.  The equality classification identifies
the extremal objects in all three languages.
\end{remark}



\section{Further directions}
\label{sec:discussion}

For the standalone parabolic problem, the present theorem is attained for
every even \(2\le k\le p-1\), whereas Section~\ref{sec:odd-refinements} shows
that sparse odd support can be much more rigid.  The natural next problem is
to determine the optimal fixed-\(k\) zero count for general odd support and
to understand the remaining early-jet configurations.  The parity-sensitive theorem leaves two natural extremal questions: determine the exact fixed-support maximum for odd \(k\ge7\), and characterize the even-support sums whose large zero sets contain an affine zero line.

Two features of the proof are genuinely prime-dimensional.  First, the
minimal-support reduction ultimately rests on the sharp Fourier uncertainty
principle on $\F_p$.  Second, the jet argument uses the total ramification of
$p$ in $\Q(\omega)$ and only the range below the critical order $p-1$.
Prime-power fields therefore present a qualitatively different problem:
subfield structure creates additional Fourier extremizers, while the local
cyclotomic model must be replaced by a higher-dimensional additive
structure.  Determining the correct analogue of the parabolic zero theorem
in that setting is a natural next question.

The cyclotomic-jet mechanism also suggests a broader problem: determine
which algebraic frequency varieties admit comparably strong control of
successive jet components and hence ambient-size-independent zero bounds.
We do not pursue that general theory here; the present paper isolates the
quadratic case, where transverse linearity and two-point fibers combine to
produce the sharp parity phenomenon.

\paragraph{Declaration of AI-assisted tools.}
OpenAI's ChatGPT was used as an auxiliary tool for proof checking,
computations, and language editing.  The author takes full responsibility for
the mathematical content and conclusions of the manuscript.

\begin{thebibliography}{99}

\bibitem{SegreOval}
B.~Segre,
\emph{Ovals in a finite projective plane},
Canad. J. Math. \textbf{7} (1955), 414--416.
\href{https://doi.org/10.4153/CJM-1955-045-X}
{doi:10.4153/CJM-1955-045-X}.



\bibitem{DonohoStark}
D.~L. Donoho and P.~B. Stark,
\emph{Uncertainty principles and signal recovery},
SIAM J. Appl. Math. \textbf{49} (1989), 906--931.
\href{https://doi.org/10.1137/0149053}{doi:10.1137/0149053}.

\bibitem{Meshulam}
R. Meshulam,
\emph{An uncertainty inequality for finite abelian groups},
European J. Combin. \textbf{27} (2006), 63--67.
\href{https://doi.org/10.1016/j.ejc.2004.07.009}{doi:10.1016/j.ejc.2004.07.009}.

\bibitem{GarciaKaraaliKatz}
S.~R. Garcia, G. Karaali and D.~J. Katz,
\emph{An improved uncertainty principle for functions with symmetry},
J. Algebra \textbf{586} (2021), 899--934.
\href{https://doi.org/10.1016/j.jalgebra.2021.07.017}{doi:10.1016/j.jalgebra.2021.07.017}.

\bibitem{DeBievre}
S. De Bi\`evre,
\emph{Complete incompatibility, support uncertainty, and
Kirkwood--Dirac nonclassicality},
Phys. Rev. Lett. \textbf{127} (2021), 190404.
\href{https://doi.org/10.1103/PhysRevLett.127.190404}{doi:10.1103/PhysRevLett.127.190404}.

\bibitem{FultonCurves}
W. Fulton,
\emph{Algebraic Curves: An Introduction to Algebraic Geometry},
Addison--Wesley, Redwood City, CA, 1989.

\bibitem{Washington}
L.~C. Washington,
\emph{Introduction to Cyclotomic Fields}, 2nd ed.,
Graduate Texts in Mathematics 83, Springer, New York, 1997.

\bibitem{WF}
W.~K. Wootters and B.~D. Fields,
\emph{Optimal state-determination by mutually unbiased measurements},
Ann. Phys. \textbf{191} (1989), 363--381.
\href{https://doi.org/10.1016/0003-4916(89)90322-9}{doi:10.1016/0003-4916(89)90322-9}.

\bibitem{Tao}
T. Tao,
\emph{An uncertainty principle for cyclic groups of prime order},
Math. Res. Lett. \textbf{12} (2005), 121--127.
\href{https://doi.org/10.4310/MRL.2005.v12.n1.a11}{doi:10.4310/MRL.2005.v12.n1.a11}.

\bibitem{FW}
V. Fiorentino and S. Weigert,
\emph{Uncertainty relations for the support of quantum states},
J. Phys. A: Math. Theor. \textbf{55} (2022), 495305.
\href{https://doi.org/10.1088/1751-8121/aca9a4}{doi:10.1088/1751-8121/aca9a4}.

\bibitem{BiroLev}
A. Bir\'o and V.~F. Lev,
\emph{Uncertainty in finite planes},
J. Funct. Anal. \textbf{281} (2021), 109026.
\href{https://doi.org/10.1016/j.jfa.2021.109026}{doi:10.1016/j.jfa.2021.109026}.

\bibitem{Lewko}
M. Lewko,
\emph{New restriction estimates for the 3-d paraboloid over finite fields},
Adv. Math. \textbf{270} (2015), 457--479.
\href{https://doi.org/10.1016/j.aim.2014.11.008}{doi:10.1016/j.aim.2014.11.008}.

\bibitem{IKL}
A. Iosevich, D. Koh and M. Lewko,
\emph{Finite field restriction estimates for the paraboloid in high even dimensions},
J. Funct. Anal. \textbf{278} (2020), 108450.
\href{https://doi.org/10.1016/j.jfa.2019.108450}{doi:10.1016/j.jfa.2019.108450}.



\bibitem{BonamiGhobber}
A. Bonami and S. Ghobber,
\emph{Equality cases for the uncertainty principle in finite Abelian groups},
Acta Sci. Math. (Szeged) \textbf{79} (2013), 507--528.

\bibitem{NicolaSTFT}
F. Nicola,
\emph{The uncertainty principle for the short-time Fourier transform on finite cyclic groups: cases of equality},
J. Funct. Anal. \textbf{284} (2023), 109924.
\href{https://doi.org/10.1016/j.jfa.2023.109924}{doi:10.1016/j.jfa.2023.109924}.

\bibitem{DeBievre2023}
S. De Bi\`evre,
\emph{Relating incompatibility, noncommutativity, uncertainty, and Kirkwood--Dirac nonclassicality},
J. Math. Phys. \textbf{64} (2023), 022202.
\href{https://doi.org/10.1063/5.0110267}{doi:10.1063/5.0110267}.

\bibitem{XuIncompatibility}
J. Xu,
\emph{Classification of incompatibility for two orthonormal bases},
Phys. Rev. A \textbf{106} (2022), 022217.
\href{https://doi.org/10.1103/PhysRevA.106.022217}{doi:10.1103/PhysRevA.106.022217}.

\bibitem{MockenhauptTao}
G. Mockenhaupt and T. Tao,
\emph{Restriction and Kakeya phenomena for finite fields},
Duke Math. J. \textbf{121} (2004), no.~1, 35--74.
\href{https://doi.org/10.1215/S0012-7094-04-12112-8}{doi:10.1215/S0012-7094-04-12112-8}.

\bibitem{LewkoKakeya}
M. Lewko,
\emph{Finite field restriction estimates based on Kakeya maximal operator estimates},
J. Eur. Math. Soc. \textbf{21} (2019), no.~12, 3649--3707.
\href{https://doi.org/10.4171/JEMS/910}{doi:10.4171/JEMS/910}.

\bibitem{IosevichMayeli}
A. Iosevich and A. Mayeli,
\emph{Uncertainty principles, restriction, Bourgain's $\Lambda_q$ theorem, and signal recovery},
Appl. Comput. Harmon. Anal. \textbf{76} (2025), 101734.
\href{https://doi.org/10.1016/j.acha.2024.101734}{doi:10.1016/j.acha.2024.101734}.

\bibitem{IosevichJamingMayeli}
A. Iosevich, P. Jaming and A. Mayeli,
\emph{Uncertainty principle, annihilating pairs and Fourier restriction},
Acta Math. Hungar. \textbf{178} (2026), 220--238.
\href{https://doi.org/10.1007/s10474-026-01589-9}{doi:10.1007/s10474-026-01589-9}.

\bibitem{GonzalezOliveira}
C. Gonz\'alez-Riquelme and D. Oliveira e Silva,
\emph{Sharp extension inequalities on finite fields},
arXiv:2405.16647 (2024), submitted.

\bibitem{BiswasEtAl}
C. Biswas, E. Carneiro, T.~C. Flock, D. Oliveira e Silva, B. Stovall and J. Tautges,
\emph{Sharp endpoint extension inequalities for the moment curve on finite fields},
Int. Math. Res. Not. IMRN \textbf{2026} (2026), no.~2, rnaf379.
\href{https://doi.org/10.1093/imrn/rnaf379}{doi:10.1093/imrn/rnaf379}.

\bibitem{HedenmalmMontes}
H. Hedenmalm and A. Montes-Rodr\'iguez,
\emph{Heisenberg uniqueness pairs and the Klein--Gordon equation},
Ann. of Math. (2) \textbf{173} (2011), no.~3, 1507--1527.

\bibitem{GiriSrivastava}
D.~K. Giri and R.~K. Srivastava,
\emph{Heisenberg uniqueness pairs for some algebraic curves in the plane},
Adv. Math. \textbf{310} (2017), 993--1016.
\href{https://doi.org/10.1016/j.aim.2017.02.019}{doi:10.1016/j.aim.2017.02.019}.

\bibitem{AlexeevCahillMixon}
B. Alexeev, J. Cahill and D.~G. Mixon,
\emph{Full spark frames},
J. Fourier Anal. Appl. \textbf{18} (2012), no.~6, 1167--1194.
\href{https://doi.org/10.1007/s00041-012-9235-4}{doi:10.1007/s00041-012-9235-4}.

\bibitem{Blanchfield}
K. Blanchfield,
\emph{Orbits of mutually unbiased bases},
J. Phys. A: Math. Theor. \textbf{47} (2014), 135303.
\href{https://doi.org/10.1088/1751-8113/47/13/135303}{doi:10.1088/1751-8113/47/13/135303}.

\bibitem{JafarpourDuarteCalderbank}
S. Jafarpour, M.~F. Duarte and R. Calderbank,
\emph{Beyond worst-case reconstruction in deterministic compressed sensing},
in: 2012 IEEE International Symposium on Information Theory Proceedings,
2012, 1852--1856.
\href{https://doi.org/10.1109/ISIT.2012.6283601}{doi:10.1109/ISIT.2012.6283601}.



\bibitem{SmithGroups}
K.~T. Smith,
\emph{The uncertainty principle on groups},
SIAM J. Appl. Math. \textbf{50} (1990), no.~3, 876--882.
\href{https://doi.org/10.1137/0150051}{doi:10.1137/0150051}.

\bibitem{WigdersonSurvey}
A. Wigderson and Y. Wigderson,
\emph{The uncertainty principle: variations on a theme},
Bull. Amer. Math. Soc. (N.S.) \textbf{58} (2021), no.~2, 225--261.
\href{https://doi.org/10.1090/bull/1715}{doi:10.1090/bull/1715}.

\bibitem{EvansIsaacs}
R.~J. Evans and I.~M. Isaacs,
\emph{Generalized Vandermonde determinants and roots of unity of prime order},
Proc. Amer. Math. Soc. \textbf{58} (1976), 51--54.

\bibitem{LoukakiChebotarev}
M. Loukaki,
\emph{Chebotarev's theorem for cyclic groups of order $pq$ and an uncertainty principle},
Bull. Lond. Math. Soc. \textbf{57} (2025), 3841--3856.
\href{https://doi.org/10.1112/blms.70192}{doi:10.1112/blms.70192}.

\bibitem{BandyopadhyayMUB}
S. Bandyopadhyay, P.~O. Boykin, V.~P. Roychowdhury and F. Vatan,
\emph{A new proof for the existence of mutually unbiased bases},
Algorithmica \textbf{34} (2002), no.~4, 512--528.
\href{https://doi.org/10.1007/s00453-002-0980-7}{doi:10.1007/s00453-002-0980-7}.

\bibitem{KlappeneckerRottelerConstruction}
A. Klappenecker and M. R\"otteler,
\emph{Constructions of mutually unbiased bases},
in: G.~L. Mullen, A. Poli and H. Stichtenoth (eds.),
Finite Fields and Applications, Lecture Notes in Comput. Sci. \textbf{2948},
Springer, Berlin, 2004, 137--144.
\href{https://doi.org/10.1007/978-3-540-24633-6_10}{doi:10.1007/978-3-540-24633-6\_10}.

\bibitem{DurtMUB}
T. Durt, B.-G. Englert, I. Bengtsson and K. \ifmmode\dot{Z}\else\.Z\fi yczkowski,
\emph{On mutually unbiased bases},
Int. J. Quantum Inf. \textbf{8} (2010), no.~4, 535--640.
\href{https://doi.org/10.1142/S0219749910006502}{doi:10.1142/S0219749910006502}.

\bibitem{McNultyWeigertReview}
D. McNulty and S. Weigert,
\emph{Mutually unbiased bases in composite dimensions---a review},
Quantum \textbf{10} (2026), 2051.
\href{https://doi.org/10.22331/q-2026-04-01-2051}{doi:10.22331/q-2026-04-01-2051}.

\bibitem{KlappeneckerRottelerDesign}
A. Klappenecker and M. R\"otteler,
\emph{Mutually unbiased bases are complex projective $2$-designs},
in: Proc. IEEE Int. Symp. Inform. Theory (ISIT 2005), 1740--1744.
\href{https://doi.org/10.1109/ISIT.2005.1523643}{doi:10.1109/ISIT.2005.1523643}.

\bibitem{RoyScottDesign}
A. Roy and A.~J. Scott,
\emph{Weighted complex projective $2$-designs from bases: optimal state determination by orthogonal measurements},
J. Math. Phys. \textbf{48} (2007), 072110.
\href{https://doi.org/10.1063/1.2748617}{doi:10.1063/1.2748617}.

\bibitem{EladBruckstein}
M. Elad and A.~M. Bruckstein,
\emph{A generalized uncertainty principle and sparse representation in pairs of bases},
IEEE Trans. Inform. Theory \textbf{48} (2002), no.~9, 2558--2567.
\href{https://doi.org/10.1109/TIT.2002.801410}{doi:10.1109/TIT.2002.801410}.

\bibitem{CasazzaKutyniok}
P.~G. Casazza and G. Kutyniok (eds.),
\emph{Finite Frames: Theory and Applications},
Applied and Numerical Harmonic Analysis, Birkh\"auser, Boston, 2013.
\href{https://doi.org/10.1007/978-0-8176-8373-3}{doi:10.1007/978-0-8176-8373-3}.

\bibitem{BodmannPaulsen}
B.~G. Bodmann and V.~I. Paulsen,
\emph{Frames, graphs and erasures},
Linear Algebra Appl. \textbf{404} (2005), 118--146.
\href{https://doi.org/10.1016/j.laa.2005.02.016}{doi:10.1016/j.laa.2005.02.016}.

\bibitem{LiFrameErasures}
D. Li, J. Leng, T. Huang and Q. Gao,
\emph{Frame expansions with probabilistic erasures},
Digit. Signal Process. \textbf{72} (2018), 75--82.
\href{https://doi.org/10.1016/j.dsp.2017.10.002}{doi:10.1016/j.dsp.2017.10.002}.

\bibitem{LiTNS2026}
D. Li,
\emph{Robustness of frames and totally nonsingular matrices},
SIAM J. Matrix Anal. Appl. \textbf{47} (2026), no.~1, 412--428.
\href{https://doi.org/10.1137/25M1766310}{doi:10.1137/25M1766310}.

\bibitem{HuffmanPless}
W.~C. Huffman and V. Pless,
\emph{Fundamentals of Error-Correcting Codes},
Cambridge University Press, Cambridge, 2003.
\href{https://doi.org/10.1017/CBO9780511807077}{doi:10.1017/CBO9780511807077}.

\bibitem{Oxley}
J. Oxley,
\emph{Matroid Theory}, 2nd ed.,
Oxford University Press, Oxford, 2011.
\href{https://doi.org/10.1093/acprof:oso/9780198566946.001.0001}{doi:10.1093/acprof:oso/9780198566946.001.0001}.

\end{thebibliography}
\end{document}
